% Options for packages loaded elsewhere
% Options for packages loaded elsewhere
\PassOptionsToPackage{unicode}{hyperref}
\PassOptionsToPackage{hyphens}{url}
\PassOptionsToPackage{dvipsnames,svgnames,x11names}{xcolor}
%
\documentclass[
  letterpaper,
  DIV=11,
  numbers=noendperiod]{scrartcl}
\usepackage{xcolor}
\usepackage{amsmath,amssymb}
\setcounter{secnumdepth}{-\maxdimen} % remove section numbering
\usepackage{iftex}
\ifPDFTeX
  \usepackage[T1]{fontenc}
  \usepackage[utf8]{inputenc}
  \usepackage{textcomp} % provide euro and other symbols
\else % if luatex or xetex
  \usepackage{unicode-math} % this also loads fontspec
  \defaultfontfeatures{Scale=MatchLowercase}
  \defaultfontfeatures[\rmfamily]{Ligatures=TeX,Scale=1}
\fi
\usepackage{lmodern}
\ifPDFTeX\else
  % xetex/luatex font selection
\fi
% Use upquote if available, for straight quotes in verbatim environments
\IfFileExists{upquote.sty}{\usepackage{upquote}}{}
\IfFileExists{microtype.sty}{% use microtype if available
  \usepackage[]{microtype}
  \UseMicrotypeSet[protrusion]{basicmath} % disable protrusion for tt fonts
}{}
\makeatletter
\@ifundefined{KOMAClassName}{% if non-KOMA class
  \IfFileExists{parskip.sty}{%
    \usepackage{parskip}
  }{% else
    \setlength{\parindent}{0pt}
    \setlength{\parskip}{6pt plus 2pt minus 1pt}}
}{% if KOMA class
  \KOMAoptions{parskip=half}}
\makeatother
% Make \paragraph and \subparagraph free-standing
\makeatletter
\ifx\paragraph\undefined\else
  \let\oldparagraph\paragraph
  \renewcommand{\paragraph}{
    \@ifstar
      \xxxParagraphStar
      \xxxParagraphNoStar
  }
  \newcommand{\xxxParagraphStar}[1]{\oldparagraph*{#1}\mbox{}}
  \newcommand{\xxxParagraphNoStar}[1]{\oldparagraph{#1}\mbox{}}
\fi
\ifx\subparagraph\undefined\else
  \let\oldsubparagraph\subparagraph
  \renewcommand{\subparagraph}{
    \@ifstar
      \xxxSubParagraphStar
      \xxxSubParagraphNoStar
  }
  \newcommand{\xxxSubParagraphStar}[1]{\oldsubparagraph*{#1}\mbox{}}
  \newcommand{\xxxSubParagraphNoStar}[1]{\oldsubparagraph{#1}\mbox{}}
\fi
\makeatother

\usepackage{color}
\usepackage{fancyvrb}
\newcommand{\VerbBar}{|}
\newcommand{\VERB}{\Verb[commandchars=\\\{\}]}
\DefineVerbatimEnvironment{Highlighting}{Verbatim}{commandchars=\\\{\}}
% Add ',fontsize=\small' for more characters per line
\usepackage{framed}
\definecolor{shadecolor}{RGB}{241,243,245}
\newenvironment{Shaded}{\begin{snugshade}}{\end{snugshade}}
\newcommand{\AlertTok}[1]{\textcolor[rgb]{0.68,0.00,0.00}{#1}}
\newcommand{\AnnotationTok}[1]{\textcolor[rgb]{0.37,0.37,0.37}{#1}}
\newcommand{\AttributeTok}[1]{\textcolor[rgb]{0.40,0.45,0.13}{#1}}
\newcommand{\BaseNTok}[1]{\textcolor[rgb]{0.68,0.00,0.00}{#1}}
\newcommand{\BuiltInTok}[1]{\textcolor[rgb]{0.00,0.23,0.31}{#1}}
\newcommand{\CharTok}[1]{\textcolor[rgb]{0.13,0.47,0.30}{#1}}
\newcommand{\CommentTok}[1]{\textcolor[rgb]{0.37,0.37,0.37}{#1}}
\newcommand{\CommentVarTok}[1]{\textcolor[rgb]{0.37,0.37,0.37}{\textit{#1}}}
\newcommand{\ConstantTok}[1]{\textcolor[rgb]{0.56,0.35,0.01}{#1}}
\newcommand{\ControlFlowTok}[1]{\textcolor[rgb]{0.00,0.23,0.31}{\textbf{#1}}}
\newcommand{\DataTypeTok}[1]{\textcolor[rgb]{0.68,0.00,0.00}{#1}}
\newcommand{\DecValTok}[1]{\textcolor[rgb]{0.68,0.00,0.00}{#1}}
\newcommand{\DocumentationTok}[1]{\textcolor[rgb]{0.37,0.37,0.37}{\textit{#1}}}
\newcommand{\ErrorTok}[1]{\textcolor[rgb]{0.68,0.00,0.00}{#1}}
\newcommand{\ExtensionTok}[1]{\textcolor[rgb]{0.00,0.23,0.31}{#1}}
\newcommand{\FloatTok}[1]{\textcolor[rgb]{0.68,0.00,0.00}{#1}}
\newcommand{\FunctionTok}[1]{\textcolor[rgb]{0.28,0.35,0.67}{#1}}
\newcommand{\ImportTok}[1]{\textcolor[rgb]{0.00,0.46,0.62}{#1}}
\newcommand{\InformationTok}[1]{\textcolor[rgb]{0.37,0.37,0.37}{#1}}
\newcommand{\KeywordTok}[1]{\textcolor[rgb]{0.00,0.23,0.31}{\textbf{#1}}}
\newcommand{\NormalTok}[1]{\textcolor[rgb]{0.00,0.23,0.31}{#1}}
\newcommand{\OperatorTok}[1]{\textcolor[rgb]{0.37,0.37,0.37}{#1}}
\newcommand{\OtherTok}[1]{\textcolor[rgb]{0.00,0.23,0.31}{#1}}
\newcommand{\PreprocessorTok}[1]{\textcolor[rgb]{0.68,0.00,0.00}{#1}}
\newcommand{\RegionMarkerTok}[1]{\textcolor[rgb]{0.00,0.23,0.31}{#1}}
\newcommand{\SpecialCharTok}[1]{\textcolor[rgb]{0.37,0.37,0.37}{#1}}
\newcommand{\SpecialStringTok}[1]{\textcolor[rgb]{0.13,0.47,0.30}{#1}}
\newcommand{\StringTok}[1]{\textcolor[rgb]{0.13,0.47,0.30}{#1}}
\newcommand{\VariableTok}[1]{\textcolor[rgb]{0.07,0.07,0.07}{#1}}
\newcommand{\VerbatimStringTok}[1]{\textcolor[rgb]{0.13,0.47,0.30}{#1}}
\newcommand{\WarningTok}[1]{\textcolor[rgb]{0.37,0.37,0.37}{\textit{#1}}}

\usepackage{longtable,booktabs,array}
\usepackage{calc} % for calculating minipage widths
% Correct order of tables after \paragraph or \subparagraph
\usepackage{etoolbox}
\makeatletter
\patchcmd\longtable{\par}{\if@noskipsec\mbox{}\fi\par}{}{}
\makeatother
% Allow footnotes in longtable head/foot
\IfFileExists{footnotehyper.sty}{\usepackage{footnotehyper}}{\usepackage{footnote}}
\makesavenoteenv{longtable}
\usepackage{graphicx}
\makeatletter
\newsavebox\pandoc@box
\newcommand*\pandocbounded[1]{% scales image to fit in text height/width
  \sbox\pandoc@box{#1}%
  \Gscale@div\@tempa{\textheight}{\dimexpr\ht\pandoc@box+\dp\pandoc@box\relax}%
  \Gscale@div\@tempb{\linewidth}{\wd\pandoc@box}%
  \ifdim\@tempb\p@<\@tempa\p@\let\@tempa\@tempb\fi% select the smaller of both
  \ifdim\@tempa\p@<\p@\scalebox{\@tempa}{\usebox\pandoc@box}%
  \else\usebox{\pandoc@box}%
  \fi%
}
% Set default figure placement to htbp
\def\fps@figure{htbp}
\makeatother





\setlength{\emergencystretch}{3em} % prevent overfull lines

\providecommand{\tightlist}{%
  \setlength{\itemsep}{0pt}\setlength{\parskip}{0pt}}



 


\KOMAoption{captions}{tableheading}
\makeatletter
\@ifpackageloaded{tcolorbox}{}{\usepackage[skins,breakable]{tcolorbox}}
\@ifpackageloaded{fontawesome5}{}{\usepackage{fontawesome5}}
\definecolor{quarto-callout-color}{HTML}{909090}
\definecolor{quarto-callout-note-color}{HTML}{0758E5}
\definecolor{quarto-callout-important-color}{HTML}{CC1914}
\definecolor{quarto-callout-warning-color}{HTML}{EB9113}
\definecolor{quarto-callout-tip-color}{HTML}{00A047}
\definecolor{quarto-callout-caution-color}{HTML}{FC5300}
\definecolor{quarto-callout-color-frame}{HTML}{acacac}
\definecolor{quarto-callout-note-color-frame}{HTML}{4582ec}
\definecolor{quarto-callout-important-color-frame}{HTML}{d9534f}
\definecolor{quarto-callout-warning-color-frame}{HTML}{f0ad4e}
\definecolor{quarto-callout-tip-color-frame}{HTML}{02b875}
\definecolor{quarto-callout-caution-color-frame}{HTML}{fd7e14}
\makeatother
\makeatletter
\@ifpackageloaded{caption}{}{\usepackage{caption}}
\AtBeginDocument{%
\ifdefined\contentsname
  \renewcommand*\contentsname{Table of contents}
\else
  \newcommand\contentsname{Table of contents}
\fi
\ifdefined\listfigurename
  \renewcommand*\listfigurename{List of Figures}
\else
  \newcommand\listfigurename{List of Figures}
\fi
\ifdefined\listtablename
  \renewcommand*\listtablename{List of Tables}
\else
  \newcommand\listtablename{List of Tables}
\fi
\ifdefined\figurename
  \renewcommand*\figurename{Figure}
\else
  \newcommand\figurename{Figure}
\fi
\ifdefined\tablename
  \renewcommand*\tablename{Table}
\else
  \newcommand\tablename{Table}
\fi
}
\@ifpackageloaded{float}{}{\usepackage{float}}
\floatstyle{ruled}
\@ifundefined{c@chapter}{\newfloat{codelisting}{h}{lop}}{\newfloat{codelisting}{h}{lop}[chapter]}
\floatname{codelisting}{Listing}
\newcommand*\listoflistings{\listof{codelisting}{List of Listings}}
\makeatother
\makeatletter
\makeatother
\makeatletter
\@ifpackageloaded{caption}{}{\usepackage{caption}}
\@ifpackageloaded{subcaption}{}{\usepackage{subcaption}}
\makeatother
\usepackage{bookmark}
\IfFileExists{xurl.sty}{\usepackage{xurl}}{} % add URL line breaks if available
\urlstyle{same}
\hypersetup{
  pdftitle={Data Manipulation Challenge},
  colorlinks=true,
  linkcolor={blue},
  filecolor={Maroon},
  citecolor={Blue},
  urlcolor={Blue},
  pdfcreator={LaTeX via pandoc}}


\title{Data Manipulation Challenge}
\usepackage{etoolbox}
\makeatletter
\providecommand{\subtitle}[1]{% add subtitle to \maketitle
  \apptocmd{\@title}{\par {\large #1 \par}}{}{}
}
\makeatother
\subtitle{A Mental Model for Method Chaining in Pandas}
\author{}
\date{}
\begin{document}
\maketitle


\section{🔗 Data Manipulation Challenge - A Mental Model for Method
Chaining in
Pandas}\label{data-manipulation-challenge---a-mental-model-for-method-chaining-in-pandas}

\begin{tcolorbox}[enhanced jigsaw, leftrule=.75mm, bottomtitle=1mm, toprule=.15mm, toptitle=1mm, left=2mm, colbacktitle=quarto-callout-important-color!10!white, rightrule=.15mm, titlerule=0mm, colback=white, colframe=quarto-callout-important-color-frame, title=\textcolor{quarto-callout-important-color}{\faExclamation}\hspace{0.5em}{📊 Challenge Requirements In Section
\hyperref[student-analysis-section]{Student Analysis Section}}, arc=.35mm, coltitle=black, opacitybacktitle=0.6, breakable, opacityback=0, bottomrule=.15mm]

\begin{itemize}
\tightlist
\item
  Complete all discussion questions for the seven mental models (plus
  some extra requirements for higher grades)
\end{itemize}

\end{tcolorbox}

\begin{tcolorbox}[enhanced jigsaw, leftrule=.75mm, bottomtitle=1mm, toprule=.15mm, toptitle=1mm, left=2mm, colbacktitle=quarto-callout-important-color!10!white, rightrule=.15mm, titlerule=0mm, colback=white, colframe=quarto-callout-important-color-frame, title=\textcolor{quarto-callout-important-color}{\faExclamation}\hspace{0.5em}{🎯 Note on Python Usage}, arc=.35mm, coltitle=black, opacitybacktitle=0.6, breakable, opacityback=0, bottomrule=.15mm]

\textbf{Recommended Workflow: Use Your Existing Virtual Environment} If
you completed the Tech Setup Challenge Part 2, you already have a
virtual environment set up! Here's how to use it for this new challenge:

\begin{enumerate}
\def\labelenumi{\arabic{enumi}.}
\tightlist
\item
  \textbf{Clone this new challenge repository} (see Getting Started
  section below)
\item
  \textbf{Open the cloned repository in Cursor}
\item
  \textbf{Set this project to use your existing Python interpreter:}

  \begin{itemize}
  \tightlist
  \item
    Press \texttt{Ctrl+Shift+P} → ``Python: Select Interpreter''
  \item
    Navigate to and choose the interpreter from your existing virtual
    environment (e.g.,
    \texttt{your-previous-project/venv/Scripts/python.exe})
  \end{itemize}
\item
  \textbf{Activate the environment in your terminal:}

  \begin{itemize}
  \tightlist
  \item
    Open terminal in Cursor (`Ctrl + ``)
  \item
    Navigate to your previous project folder where you have the
    \texttt{venv} folder
  \item
    \textbf{💡 Pro tip:} You can quickly navigate by typing \texttt{cd}
    followed by dragging the folder from your file explorer into the
    terminal
  \item
    Activate using the appropriate command for your system:

    \begin{itemize}
    \tightlist
    \item
      \textbf{Windows Command Prompt:}
      \texttt{venv\textbackslash{}Scripts\textbackslash{}activate}
    \item
      \textbf{Windows PowerShell:}
      \texttt{.\textbackslash{}venv\textbackslash{}Scripts\textbackslash{}Activate.ps1}
    \item
      \textbf{Mac/Linux:} \texttt{source\ venv/bin/activate}
    \end{itemize}
  \item
    You should see \texttt{(venv)} at the beginning of your terminal
    prompt
  \end{itemize}
\item
  \textbf{Install additional packages if needed:}
  \texttt{pip\ install\ pandas\ numpy\ matplotlib\ seaborn}
\end{enumerate}

\begin{quote}
\textbf{⚠️ Cloud Storage Warning}

\textbf{Avoid using Google Drive, OneDrive, or other cloud storage for
Python projects!} These services can cause issues with: - Package
installations failing due to file locking - Virtual environment
corruption - Slow performance during pip operations

\textbf{Best practice:} Keep your Python projects in a local folder like
\texttt{C:\textbackslash{}Users\textbackslash{}YourName\textbackslash{}Documents\textbackslash{}}
or \texttt{\textasciitilde{}/Documents/} instead of cloud-synced
folders.
\end{quote}

\textbf{Alternative: Create a New Virtual Environment} If you prefer a
fresh environment, follow the Quarto documentation:
\url{https://quarto.org/docs/projects/virtual-environments.html}. Be
sure to follow the instructions to activate the environment, set it up
as your default Python interpreter for the project, and install the
necessary packages (e.g.~pandas) for this challenge. For installing the
packages, you can use the \texttt{pip\ install\ -r\ requirements.txt}
command since you already have the requirements.txt file in your
project. Some steps do take a bit of time, so be patient.

\textbf{Why This Works:} Virtual environments are portable - you can use
the same environment across multiple projects, and Cursor automatically
activates it when you select the interpreter!

\end{tcolorbox}

\subsection{The Problem: Mastering Data Manipulation Through Method
Chaining}\label{the-problem-mastering-data-manipulation-through-method-chaining}

\textbf{Core Question:} How can we efficiently manipulate datasets using
\texttt{pandas} method chaining to answer complex business questions?

\textbf{The Challenge:} Real-world data analysis requires combining
multiple data manipulation techniques in sequence. Rather than creating
intermediate variables at each step, method chaining allows us to write
clean, readable code that flows logically from one operation to the
next.

\textbf{Our Approach:} We'll work with ZappTech's shipment data to
answer critical business questions about service levels and
cross-category orders, using the seven mental models of data
manipulation through pandas method chaining.

\begin{tcolorbox}[enhanced jigsaw, leftrule=.75mm, bottomtitle=1mm, toprule=.15mm, toptitle=1mm, left=2mm, colbacktitle=quarto-callout-warning-color!10!white, rightrule=.15mm, titlerule=0mm, colback=white, colframe=quarto-callout-warning-color-frame, title=\textcolor{quarto-callout-warning-color}{\faExclamationTriangle}\hspace{0.5em}{⚠️ AI Partnership Required}, arc=.35mm, coltitle=black, opacitybacktitle=0.6, breakable, opacityback=0, bottomrule=.15mm]

This challenge pushes boundaries intentionally. You'll tackle problems
that normally require weeks of study, but with Cursor AI as your partner
(and your brain keeping it honest), you can accomplish more than you
thought possible.

\textbf{The new reality:} The four stages of competence are Ignorance →
Awareness → Learning → Mastery. AI lets us produce Mastery-level work
while operating primarily in the Awareness stage. I focus on awareness
training, you leverage AI for execution, and together we create outputs
that used to require years of dedicated study.

\end{tcolorbox}

\subsection{The Seven Mental Models of Data
Manipulation}\label{the-seven-mental-models-of-data-manipulation}

The seven most important ways we manipulate datasets are:

\begin{enumerate}
\def\labelenumi{\arabic{enumi}.}
\tightlist
\item
  \textbf{Assign:} Add new variables with calculations and
  transformations
\item
  \textbf{Subset:} Filter data based on conditions or select specific
  columns
\item
  \textbf{Drop:} Remove unwanted variables or observations
\item
  \textbf{Sort:} Arrange data by values or indices
\item
  \textbf{Aggregate:} Summarize data using functions like mean, sum,
  count
\item
  \textbf{Merge:} Combine information from multiple datasets
\item
  \textbf{Split-Apply-Combine:} Group data and apply functions within
  groups
\end{enumerate}

\subsection{Data and Business Context}\label{data-and-business-context}

We analyze ZappTech's shipment data, which contains information about
product deliveries across multiple categories. This dataset is ideal for
our analysis because:

\begin{itemize}
\tightlist
\item
  \textbf{Real Business Questions:} CEO wants to understand service
  levels and cross-category shopping patterns
\item
  \textbf{Multiple Data Sources:} Requires merging shipment data with
  product category information
\item
  \textbf{Complex Relationships:} Service levels may vary by product
  category, and customers may order across categories
\item
  \textbf{Method Chaining Practice:} Perfect for demonstrating all seven
  mental models in sequence
\end{itemize}

\subsection{Data Loading and Initial
Exploration}\label{data-loading-and-initial-exploration}

Let's start by loading the ZappTech shipment data and understanding what
we're working with.

\phantomsection\label{load-data}
\begin{Shaded}
\begin{Highlighting}[]
\ImportTok{import}\NormalTok{ pandas }\ImportTok{as}\NormalTok{ pd}
\ImportTok{import}\NormalTok{ numpy }\ImportTok{as}\NormalTok{ np}
\ImportTok{import}\NormalTok{ matplotlib.pyplot }\ImportTok{as}\NormalTok{ plt}
\ImportTok{import}\NormalTok{ seaborn }\ImportTok{as}\NormalTok{ sns}
\ImportTok{from}\NormalTok{ datetime }\ImportTok{import}\NormalTok{ datetime, timedelta}

\CommentTok{\# Load the shipment data}
\NormalTok{shipments\_df }\OperatorTok{=}\NormalTok{ pd.read\_csv(}
    \StringTok{"https://raw.githubusercontent.com/flyaflya/persuasive/main/shipments.csv"}\NormalTok{, }
\NormalTok{    parse\_dates}\OperatorTok{=}\NormalTok{[}\StringTok{\textquotesingle{}plannedShipDate\textquotesingle{}}\NormalTok{, }\StringTok{\textquotesingle{}actualShipDate\textquotesingle{}}\NormalTok{]}
\NormalTok{)}

\CommentTok{\# Load product line data}
\NormalTok{product\_line\_df }\OperatorTok{=}\NormalTok{ pd.read\_csv(}
    \StringTok{"https://raw.githubusercontent.com/flyaflya/persuasive/main/productLine.csv"}
\NormalTok{)}

\CommentTok{\# Reduce dataset size for faster processing (4,000 rows instead of 96,805 rows)}
\NormalTok{shipments\_df }\OperatorTok{=}\NormalTok{ shipments\_df.head(}\DecValTok{4000}\NormalTok{)}

\BuiltInTok{print}\NormalTok{(}\StringTok{"Shipments data shape:"}\NormalTok{, shipments\_df.shape)}
\BuiltInTok{print}\NormalTok{(}\StringTok{"}\CharTok{\textbackslash{}n}\StringTok{Shipments data columns:"}\NormalTok{, shipments\_df.columns.tolist())}
\BuiltInTok{print}\NormalTok{(}\StringTok{"}\CharTok{\textbackslash{}n}\StringTok{First few rows of shipments data:"}\NormalTok{)}
\BuiltInTok{print}\NormalTok{(shipments\_df.head(}\DecValTok{10}\NormalTok{))}

\BuiltInTok{print}\NormalTok{(}\StringTok{"}\CharTok{\textbackslash{}n}\StringTok{"} \OperatorTok{+} \StringTok{"="}\OperatorTok{*}\DecValTok{50}\NormalTok{)}
\BuiltInTok{print}\NormalTok{(}\StringTok{"Product line data shape:"}\NormalTok{, product\_line\_df.shape)}
\BuiltInTok{print}\NormalTok{(}\StringTok{"}\CharTok{\textbackslash{}n}\StringTok{Product line data columns:"}\NormalTok{, product\_line\_df.columns.tolist())}
\BuiltInTok{print}\NormalTok{(}\StringTok{"}\CharTok{\textbackslash{}n}\StringTok{First few rows of product line data:"}\NormalTok{)}
\BuiltInTok{print}\NormalTok{(product\_line\_df.head(}\DecValTok{10}\NormalTok{))}
\end{Highlighting}
\end{Shaded}

\begin{verbatim}
Shipments data shape: (4000, 5)

Shipments data columns: ['shipID', 'plannedShipDate', 'actualShipDate', 'partID', 'quantity']

First few rows of shipments data:
   shipID plannedShipDate actualShipDate       partID  quantity
0   10001      2013-11-06     2013-10-04  part92b16c5         6
1   10002      2013-10-15     2013-10-04   part66983b         2
2   10003      2013-10-25     2013-10-07  part8e36f25         1
3   10004      2013-10-14     2013-10-08  part30f5de0         1
4   10005      2013-10-14     2013-10-08  part9d64d35         6
5   10006      2013-10-14     2013-10-08  part6cd6167        15
6   10007      2013-10-14     2013-10-08  parta4d5fd1         2
7   10008      2013-10-14     2013-10-08  part08cadf5         1
8   10009      2013-10-14     2013-10-08  part5cc4989        10
9   10010      2013-10-14     2013-10-08  part912ae4c         1

==================================================
Product line data shape: (11997, 3)

Product line data columns: ['partID', 'productLine', 'prodCategory']

First few rows of product line data:
        partID productLine prodCategory
0  part00005ba      line4c      Liquids
1  part000b57d      line61     Machines
2  part00123bf      linec1  Marketables
3  part0021fc9      line61     Machines
4  part0027e86      line2f     Machines
5  part002ed95      line4c      Liquids
6  part0030856      lineb8     Machines
7  part0033dfd      line49      Liquids
8  part0037a2a      linea3  Marketables
9  part003caee      linea3  Marketables
\end{verbatim}

\begin{tcolorbox}[enhanced jigsaw, leftrule=.75mm, bottomtitle=1mm, toprule=.15mm, toptitle=1mm, left=2mm, colbacktitle=quarto-callout-note-color!10!white, rightrule=.15mm, titlerule=0mm, colback=white, colframe=quarto-callout-note-color-frame, title=\textcolor{quarto-callout-note-color}{\faInfo}\hspace{0.5em}{💡 Understanding the Data}, arc=.35mm, coltitle=black, opacitybacktitle=0.6, breakable, opacityback=0, bottomrule=.15mm]

\textbf{Shipments Data:} Contains individual line items for each
shipment, including: - \texttt{shipID}: Unique identifier for each
shipment - \texttt{partID}: Product identifier -
\texttt{plannedShipDate}: When the shipment was supposed to go out -
\texttt{actualShipDate}: When it actually shipped - \texttt{quantity}:
How many units were shipped

\textbf{Product Category and Line Data:} Contains product category
information: - \texttt{partID}: Links to shipments data -
\texttt{productLine}: The category each product belongs to -
\texttt{prodCategory}: The category each product belongs to

\textbf{Business Questions We'll Answer:} 1. Does service level (on-time
shipments) vary across product categories? 2. How often do orders
include products from more than one category?

\end{tcolorbox}

\subsection{The Seven Mental Models: A Progressive Learning
Journey}\label{the-seven-mental-models-a-progressive-learning-journey}

Now we'll work through each of the seven mental models using method
chaining, starting simple and building complexity.

\subsubsection{1. Assign: Adding New
Variables}\label{assign-adding-new-variables}

\textbf{Mental Model:} Create new columns with calculations and
transformations.

Let's start by calculating whether each shipment was late:

\phantomsection\label{mental-model-1-assign}
\begin{Shaded}
\begin{Highlighting}[]
\CommentTok{\# Simple assignment {-} calculate if shipment was late}
\NormalTok{shipments\_with\_lateness }\OperatorTok{=}\NormalTok{ (}
\NormalTok{    shipments\_df}
\NormalTok{    .assign(}
\NormalTok{        is\_late}\OperatorTok{=}\KeywordTok{lambda}\NormalTok{ df: df[}\StringTok{\textquotesingle{}actualShipDate\textquotesingle{}}\NormalTok{] }\OperatorTok{\textgreater{}}\NormalTok{ df[}\StringTok{\textquotesingle{}plannedShipDate\textquotesingle{}}\NormalTok{],}
\NormalTok{        days\_late}\OperatorTok{=}\KeywordTok{lambda}\NormalTok{ df: (df[}\StringTok{\textquotesingle{}actualShipDate\textquotesingle{}}\NormalTok{] }\OperatorTok{{-}}\NormalTok{ df[}\StringTok{\textquotesingle{}plannedShipDate\textquotesingle{}}\NormalTok{]).dt.days}
\NormalTok{    )}
\NormalTok{)}

\BuiltInTok{print}\NormalTok{(}\StringTok{"Added lateness calculations:"}\NormalTok{)}
\BuiltInTok{print}\NormalTok{(shipments\_with\_lateness[[}\StringTok{\textquotesingle{}shipID\textquotesingle{}}\NormalTok{, }\StringTok{\textquotesingle{}plannedShipDate\textquotesingle{}}\NormalTok{, }\StringTok{\textquotesingle{}actualShipDate\textquotesingle{}}\NormalTok{, }\StringTok{\textquotesingle{}is\_late\textquotesingle{}}\NormalTok{, }\StringTok{\textquotesingle{}days\_late\textquotesingle{}}\NormalTok{]].head())}
\end{Highlighting}
\end{Shaded}

\begin{verbatim}
Added lateness calculations:
   shipID plannedShipDate actualShipDate  is_late  days_late
0   10001      2013-11-06     2013-10-04    False        -33
1   10002      2013-10-15     2013-10-04    False        -11
2   10003      2013-10-25     2013-10-07    False        -18
3   10004      2013-10-14     2013-10-08    False         -6
4   10005      2013-10-14     2013-10-08    False         -6
\end{verbatim}

\begin{tcolorbox}[enhanced jigsaw, leftrule=.75mm, bottomtitle=1mm, toprule=.15mm, toptitle=1mm, left=2mm, colbacktitle=quarto-callout-tip-color!10!white, rightrule=.15mm, titlerule=0mm, colback=white, colframe=quarto-callout-tip-color-frame, title=\textcolor{quarto-callout-tip-color}{\faLightbulb}\hspace{0.5em}{💡 Method Chaining Tip for New Python Users}, arc=.35mm, coltitle=black, opacitybacktitle=0.6, breakable, opacityback=0, bottomrule=.15mm]

\textbf{Why use \texttt{lambda\ df:}?} When chaining methods, we need to
reference the current state of the dataframe. The \texttt{lambda\ df:}
tells pandas ``use the current dataframe in this calculation.'' Without
it, pandas would look for a variable called \texttt{df} that doesn't
exist.

\textbf{Alternative approach:} You could also write this as separate
steps, but method chaining keeps related operations together and makes
the code more readable.

\end{tcolorbox}

\begin{tcolorbox}[enhanced jigsaw, leftrule=.75mm, bottomtitle=1mm, toprule=.15mm, toptitle=1mm, left=2mm, colbacktitle=quarto-callout-important-color!10!white, rightrule=.15mm, titlerule=0mm, colback=white, colframe=quarto-callout-important-color-frame, title=\textcolor{quarto-callout-important-color}{\faExclamation}\hspace{0.5em}{🤔 Discussion Questions: Assign Mental Model}, arc=.35mm, coltitle=black, opacitybacktitle=0.6, breakable, opacityback=0, bottomrule=.15mm]

\textbf{Question 1: Data Types and Date Handling} - What is the
\texttt{dtype} of the \texttt{actualShipDate} series? How can you find
out using code? - Why is it important that both \texttt{actualShipDate}
and \texttt{plannedShipDate} have the same data type for comparison?

\textbf{Question 2: String vs Date Comparison} - Can you give an example
where comparing two dates as strings would yield unintuitive results,
e.g.~what happens if you try to compare ``04-11-2025'' and
``05-20-2024'' as strings vs as dates?

\textbf{Question 3: Debug This Code}

\begin{Shaded}
\begin{Highlighting}[]
\CommentTok{\# This code has an error {-} can you spot it?}
\NormalTok{shipments\_with\_lateness }\OperatorTok{=}\NormalTok{ (}
\NormalTok{    shipments\_df}
\NormalTok{    .assign(}
\NormalTok{        is\_late}\OperatorTok{=}\KeywordTok{lambda}\NormalTok{ df: df[}\StringTok{\textquotesingle{}actualShipDate\textquotesingle{}}\NormalTok{] }\OperatorTok{\textgreater{}}\NormalTok{ df[}\StringTok{\textquotesingle{}plannedShipDate\textquotesingle{}}\NormalTok{],}
\NormalTok{        days\_late}\OperatorTok{=}\KeywordTok{lambda}\NormalTok{ df: (df[}\StringTok{\textquotesingle{}actualShipDate\textquotesingle{}}\NormalTok{] }\OperatorTok{{-}}\NormalTok{ df[}\StringTok{\textquotesingle{}plannedShipDate\textquotesingle{}}\NormalTok{]).dt.days,}
\NormalTok{        lateStatement}\OperatorTok{=}\KeywordTok{lambda}\NormalTok{ df: df[}\StringTok{\textquotesingle{}is\_late\textquotesingle{}}\NormalTok{].}\BuiltInTok{apply}\NormalTok{(}\KeywordTok{lambda}\NormalTok{ x: }\StringTok{"Darn Shipment is Late"} \ControlFlowTok{if}\NormalTok{ x }\ControlFlowTok{else} \StringTok{"Shipment is on Time"}\NormalTok{)}
\NormalTok{    )}
\NormalTok{)}
\end{Highlighting}
\end{Shaded}

What's wrong with the \texttt{lateStatement} assignment and how would
you fix it?

\end{tcolorbox}

\paragraph{Briefly Give Answers to the Discussion Questions In This
Section}\label{briefly-give-answers-to-the-discussion-questions-in-this-section}

\begin{quote}
Replace this with your answers to the discussion questions
\end{quote}

\begin{Shaded}
\begin{Highlighting}[]
\CommentTok{\# Simple assignment {-} calculate if shipment was late}
\BuiltInTok{print}\NormalTok{(shipments\_df[}\StringTok{\textquotesingle{}actualShipDate\textquotesingle{}}\NormalTok{].dtype)}
\end{Highlighting}
\end{Shaded}

\begin{verbatim}
datetime64[ns]
\end{verbatim}

1.The dtype of actualShipDate is datetime64{[}ns{]} because when we
loaded the data, we used parse\_dates={[}`plannedShipDate',
`actualShipDate'{]} which converted these columns from strings to
datetime objects. 2.It is important that both actualShipDate and
plannedShipDate have the same data type for comparison because if they
are not the same data type, pandas will not be able to compare them.
.String vs Date Comparison Problem

When comparing dates as strings, the computer evaluates them
alphabetically character-by-character rather than chronologically. This
leads to incorrect results because string comparison doesn't understand
calendar logic.

The Problem:

String comparison: ``04-11-2025'' vs ``05-20-2024'' → False (wrong)

Date comparison: April 11, 2025 vs May 20, 2024 → True (correct)

3.Debug This Code Answer: The error occurs because
shipments\_df{[}`is\_late'{]} doesn't exist in the original dataframe -
it's being created in the same assign() call. The fix is to use the
lambda approach:

\begin{Shaded}
\begin{Highlighting}[]
\NormalTok{shipments\_with\_lateness }\OperatorTok{=}\NormalTok{ (}
\NormalTok{    shipments\_df}
\NormalTok{    .assign(}
\NormalTok{        is\_late}\OperatorTok{=}\KeywordTok{lambda}\NormalTok{ df: df[}\StringTok{\textquotesingle{}actualShipDate\textquotesingle{}}\NormalTok{] }\OperatorTok{\textgreater{}}\NormalTok{ df[}\StringTok{\textquotesingle{}plannedShipDate\textquotesingle{}}\NormalTok{],}
\NormalTok{        days\_late}\OperatorTok{=}\KeywordTok{lambda}\NormalTok{ df: (df[}\StringTok{\textquotesingle{}actualShipDate\textquotesingle{}}\NormalTok{] }\OperatorTok{{-}}\NormalTok{ df[}\StringTok{\textquotesingle{}plannedShipDate\textquotesingle{}}\NormalTok{]).dt.days,}
\NormalTok{        lateStatement}\OperatorTok{=}\KeywordTok{lambda}\NormalTok{ df: df[}\StringTok{\textquotesingle{}is\_late\textquotesingle{}}\NormalTok{].}\BuiltInTok{apply}\NormalTok{(}
            \KeywordTok{lambda}\NormalTok{ x: }\StringTok{"Darn Shipment is Late"} \ControlFlowTok{if}\NormalTok{ x }\ControlFlowTok{else} \StringTok{"Shipment is on Time"}
\NormalTok{        )}
\NormalTok{    )}
\NormalTok{)}
\end{Highlighting}
\end{Shaded}

\subsubsection{2. Subset: Querying Rows and Filtering
Columns}\label{subset-querying-rows-and-filtering-columns}

\textbf{Mental Model:} Query rows based on conditions and filter to keep
specific columns.

Let's query for only late shipments and filter to keep the columns we
need:

\phantomsection\label{mental-model-2-subset}
\begin{Shaded}
\begin{Highlighting}[]
\CommentTok{\# Query rows for late shipments and filter to keep specific columns}
\NormalTok{late\_shipments }\OperatorTok{=}\NormalTok{ (}
\NormalTok{    shipments\_with\_lateness}
\NormalTok{    .query(}\StringTok{\textquotesingle{}is\_late == True\textquotesingle{}}\NormalTok{)  }\CommentTok{\# Query rows where is\_late is True}
\NormalTok{    .}\BuiltInTok{filter}\NormalTok{([}\StringTok{\textquotesingle{}shipID\textquotesingle{}}\NormalTok{, }\StringTok{\textquotesingle{}partID\textquotesingle{}}\NormalTok{, }\StringTok{\textquotesingle{}plannedShipDate\textquotesingle{}}\NormalTok{, }\StringTok{\textquotesingle{}actualShipDate\textquotesingle{}}\NormalTok{, }\StringTok{\textquotesingle{}days\_late\textquotesingle{}}\NormalTok{])  }\CommentTok{\# Filter to keep specific columns}
\NormalTok{)}

\BuiltInTok{print}\NormalTok{(}\SpecialStringTok{f"Found }\SpecialCharTok{\{}\BuiltInTok{len}\NormalTok{(late\_shipments)}\SpecialCharTok{\}}\SpecialStringTok{ late shipments out of }\SpecialCharTok{\{}\BuiltInTok{len}\NormalTok{(shipments\_with\_lateness)}\SpecialCharTok{\}}\SpecialStringTok{ total"}\NormalTok{)}
\BuiltInTok{print}\NormalTok{(}\StringTok{"}\CharTok{\textbackslash{}n}\StringTok{Late shipments sample:"}\NormalTok{)}
\BuiltInTok{print}\NormalTok{(late\_shipments.head())}
\end{Highlighting}
\end{Shaded}

\begin{verbatim}
Found 456 late shipments out of 4000 total

Late shipments sample:
     shipID       partID plannedShipDate actualShipDate  days_late
776   10192  part0164a70      2013-10-09     2013-10-14          5
777   10192  part9259836      2013-10-09     2013-10-14          5
778   10192  part4526c73      2013-10-09     2013-10-14          5
779   10192  partbb47e81      2013-10-09     2013-10-14          5
780   10192  part008482f      2013-10-09     2013-10-14          5
\end{verbatim}

\begin{tcolorbox}[enhanced jigsaw, leftrule=.75mm, bottomtitle=1mm, toprule=.15mm, toptitle=1mm, left=2mm, colbacktitle=quarto-callout-note-color!10!white, rightrule=.15mm, titlerule=0mm, colback=white, colframe=quarto-callout-note-color-frame, title=\textcolor{quarto-callout-note-color}{\faInfo}\hspace{0.5em}{🔍 Understanding the Methods}, arc=.35mm, coltitle=black, opacitybacktitle=0.6, breakable, opacityback=0, bottomrule=.15mm]

\begin{itemize}
\tightlist
\item
  \textbf{\texttt{.query()}}: Query rows based on conditions (like SQL
  WHERE clause)
\item
  \textbf{\texttt{.filter()}}: Filter to keep specific columns by name
\item
  \textbf{Alternative}: You could use \texttt{.loc{[}{]}} for more
  complex row querying, but \texttt{.query()} is often more readable
\end{itemize}

\end{tcolorbox}

\begin{tcolorbox}[enhanced jigsaw, leftrule=.75mm, bottomtitle=1mm, toprule=.15mm, toptitle=1mm, left=2mm, colbacktitle=quarto-callout-important-color!10!white, rightrule=.15mm, titlerule=0mm, colback=white, colframe=quarto-callout-important-color-frame, title=\textcolor{quarto-callout-important-color}{\faExclamation}\hspace{0.5em}{🤔 Discussion Questions: Subset Mental Model}, arc=.35mm, coltitle=black, opacitybacktitle=0.6, breakable, opacityback=0, bottomrule=.15mm]

\textbf{Question 1: Query vs Boolean Indexing} - What's the difference
between using
\texttt{.query(\textquotesingle{}is\_late\ ==\ True\textquotesingle{})}
and
\texttt{{[}df{[}\textquotesingle{}is\_late\textquotesingle{}{]}\ ==\ True{]}}?
- Which approach is more readable and why?

\textbf{Question 2: Additional Row Querying} - Can you show an example
of using a variable like \texttt{late\_threshold} to query rows for
shipments that are at least \texttt{late\_threshold} days late,
e.g.~what if you wanted to query rows for shipments that are at least 5
days late?

\end{tcolorbox}

\paragraph{Briefly Give Answers to the Discussion Questions In This
Section}\label{briefly-give-answers-to-the-discussion-questions-in-this-section-1}

1 .query() - More readable, SQL-like syntax:

Boolean indexing - More verbose: Performance \& Memory .query() - Often
faster for large datasets because it uses pandas' optimized query engine

Boolean indexing - Creates an intermediate boolean series, which uses
more memory

Flexibility .query() - Can use variables with @ symbol:

Boolean indexing - Variables work directly: Bottom line: .query() is
generally more readable, especially in method chains, while boolean
indexing gives you more explicit control.

2..query() is more readable, especially in method chains.

\begin{Shaded}
\begin{Highlighting}[]
\CommentTok{\# .query() {-} reads like natural language}
\NormalTok{(shipments\_with\_lateness}
\NormalTok{    .query(}\StringTok{\textquotesingle{}is\_late == True\textquotesingle{}}\NormalTok{)}
\NormalTok{    .}\BuiltInTok{filter}\NormalTok{([}\StringTok{\textquotesingle{}shipID\textquotesingle{}}\NormalTok{, }\StringTok{\textquotesingle{}partID\textquotesingle{}}\NormalTok{, }\StringTok{\textquotesingle{}days\_late\textquotesingle{}}\NormalTok{])}
\NormalTok{)}
\end{Highlighting}
\end{Shaded}

\phantomsection\label{mental-model-2-subset-examples}
\begin{longtable}[]{@{}llll@{}}
\toprule\noalign{}
& shipID & partID & days\_late \\
\midrule\noalign{}
\endhead
\bottomrule\noalign{}
\endlastfoot
776 & 10192 & part0164a70 & 5 \\
777 & 10192 & part9259836 & 5 \\
778 & 10192 & part4526c73 & 5 \\
779 & 10192 & partbb47e81 & 5 \\
780 & 10192 & part008482f & 5 \\
... & ... & ... & ... \\
3896 & 10956 & part98c1c48 & 21 \\
3897 & 10956 & part82e69e9 & 21 \\
3898 & 10956 & partf23fd1e & 21 \\
3899 & 10956 & part825873c & 21 \\
3997 & 11001 & partd4952a8 & 11 \\
\end{longtable}

Boolean indexing breaks the flow:

\begin{Shaded}
\begin{Highlighting}[]
\CommentTok{\# Boolean indexing {-} more cluttered}
\NormalTok{(shipments\_with\_lateness[shipments\_with\_lateness[}\StringTok{\textquotesingle{}is\_late\textquotesingle{}}\NormalTok{] }\OperatorTok{==} \VariableTok{True}\NormalTok{]}
\NormalTok{    .}\BuiltInTok{filter}\NormalTok{([}\StringTok{\textquotesingle{}shipID\textquotesingle{}}\NormalTok{, }\StringTok{\textquotesingle{}partID\textquotesingle{}}\NormalTok{, }\StringTok{\textquotesingle{}days\_late\textquotesingle{}}\NormalTok{])}
\NormalTok{)}
\end{Highlighting}
\end{Shaded}

\begin{longtable}[]{@{}llll@{}}
\toprule\noalign{}
& shipID & partID & days\_late \\
\midrule\noalign{}
\endhead
\bottomrule\noalign{}
\endlastfoot
776 & 10192 & part0164a70 & 5 \\
777 & 10192 & part9259836 & 5 \\
778 & 10192 & part4526c73 & 5 \\
779 & 10192 & partbb47e81 & 5 \\
780 & 10192 & part008482f & 5 \\
... & ... & ... & ... \\
3896 & 10956 & part98c1c48 & 21 \\
3897 & 10956 & part82e69e9 & 21 \\
3898 & 10956 & partf23fd1e & 21 \\
3899 & 10956 & part825873c & 21 \\
3997 & 11001 & partd4952a8 & 11 \\
\end{longtable}

\begin{enumerate}
\def\labelenumi{\arabic{enumi}.}
\setcounter{enumi}{2}
\tightlist
\item
  Using .query() with variables: 5 days late:
\end{enumerate}

\begin{Shaded}
\begin{Highlighting}[]
\CommentTok{\# Define threshold}
\NormalTok{late\_threshold }\OperatorTok{=} \DecValTok{5}

\CommentTok{\# Query using @ symbol for variables}
\NormalTok{very\_late\_shipments }\OperatorTok{=}\NormalTok{ (}
\NormalTok{    shipments\_with\_lateness}
\NormalTok{    .query(}\StringTok{\textquotesingle{}days\_late \textgreater{}= @late\_threshold\textquotesingle{}}\NormalTok{)}
\NormalTok{    .}\BuiltInTok{filter}\NormalTok{([}\StringTok{\textquotesingle{}shipID\textquotesingle{}}\NormalTok{, }\StringTok{\textquotesingle{}partID\textquotesingle{}}\NormalTok{, }\StringTok{\textquotesingle{}days\_late\textquotesingle{}}\NormalTok{, }\StringTok{\textquotesingle{}is\_late\textquotesingle{}}\NormalTok{])}
\NormalTok{)}

\BuiltInTok{print}\NormalTok{(}\SpecialStringTok{f"Found }\SpecialCharTok{\{}\BuiltInTok{len}\NormalTok{(very\_late\_shipments)}\SpecialCharTok{\}}\SpecialStringTok{ shipments at least }\SpecialCharTok{\{}\NormalTok{late\_threshold}\SpecialCharTok{\}}\SpecialStringTok{ days late"}\NormalTok{)}
\BuiltInTok{print}\NormalTok{(very\_late\_shipments.head())}
\end{Highlighting}
\end{Shaded}

\begin{verbatim}
Found 186 shipments at least 5 days late
     shipID       partID  days_late  is_late
776   10192  part0164a70          5     True
777   10192  part9259836          5     True
778   10192  part4526c73          5     True
779   10192  partbb47e81          5     True
780   10192  part008482f          5     True
\end{verbatim}

Using boolean indexing with variables:

\begin{Shaded}
\begin{Highlighting}[]
\CommentTok{\# Define threshold}
\NormalTok{late\_threshold }\OperatorTok{=} \DecValTok{5}

\CommentTok{\# Boolean indexing with variables}
\NormalTok{very\_late\_shipments }\OperatorTok{=}\NormalTok{ (}
\NormalTok{    shipments\_with\_lateness[}
\NormalTok{        shipments\_with\_lateness[}\StringTok{\textquotesingle{}days\_late\textquotesingle{}}\NormalTok{] }\OperatorTok{\textgreater{}=}\NormalTok{ late\_threshold}
\NormalTok{    ]}
\NormalTok{    .}\BuiltInTok{filter}\NormalTok{([}\StringTok{\textquotesingle{}shipID\textquotesingle{}}\NormalTok{, }\StringTok{\textquotesingle{}partID\textquotesingle{}}\NormalTok{, }\StringTok{\textquotesingle{}days\_late\textquotesingle{}}\NormalTok{, }\StringTok{\textquotesingle{}is\_late\textquotesingle{}}\NormalTok{])}
\NormalTok{)}

\BuiltInTok{print}\NormalTok{(}\SpecialStringTok{f"Found }\SpecialCharTok{\{}\BuiltInTok{len}\NormalTok{(very\_late\_shipments)}\SpecialCharTok{\}}\SpecialStringTok{ shipments at least }\SpecialCharTok{\{}\NormalTok{late\_threshold}\SpecialCharTok{\}}\SpecialStringTok{ days late"}\NormalTok{)}
\BuiltInTok{print}\NormalTok{(very\_late\_shipments.head())}
\end{Highlighting}
\end{Shaded}

\begin{verbatim}
Found 186 shipments at least 5 days late
     shipID       partID  days_late  is_late
776   10192  part0164a70          5     True
777   10192  part9259836          5     True
778   10192  part4526c73          5     True
779   10192  partbb47e81          5     True
780   10192  part008482f          5     True
\end{verbatim}

Multiple thresholds example:

\begin{Shaded}
\begin{Highlighting}[]
\CommentTok{\# Multiple threshold variables}
\NormalTok{min\_days\_late }\OperatorTok{=} \DecValTok{5}
\NormalTok{max\_days\_late }\OperatorTok{=} \DecValTok{15}

\CommentTok{\# Using .query() with multiple variables}
\NormalTok{moderately\_late }\OperatorTok{=}\NormalTok{ (}
\NormalTok{    shipments\_with\_lateness}
\NormalTok{    .query(}\StringTok{\textquotesingle{}days\_late \textgreater{}= @min\_days\_late and days\_late \textless{}= @max\_days\_late\textquotesingle{}}\NormalTok{)}
\NormalTok{    .sort\_values(}\StringTok{\textquotesingle{}days\_late\textquotesingle{}}\NormalTok{, ascending}\OperatorTok{=}\VariableTok{False}\NormalTok{)}
\NormalTok{)}

\BuiltInTok{print}\NormalTok{(}\SpecialStringTok{f"Shipments }\SpecialCharTok{\{}\NormalTok{min\_days\_late}\SpecialCharTok{\}}\SpecialStringTok{{-}}\SpecialCharTok{\{}\NormalTok{max\_days\_late}\SpecialCharTok{\}}\SpecialStringTok{ days late: }\SpecialCharTok{\{}\BuiltInTok{len}\NormalTok{(moderately\_late)}\SpecialCharTok{\}}\SpecialStringTok{"}\NormalTok{)}
\end{Highlighting}
\end{Shaded}

\begin{verbatim}
Shipments 5-15 days late: 128
\end{verbatim}

Key difference in variable usage: .query(): Use @variable\_name syntax

Boolean indexing: Use variables directly This approach makes code more
flexible and reusable - you can easily change the threshold without
modifying the query logic.

\subsubsection{3. Drop: Removing Unwanted
Data}\label{drop-removing-unwanted-data}

\textbf{Mental Model:} Remove columns or rows you don't need.

Let's clean up our data by removing unnecessary columns:

\phantomsection\label{mental-model-3-drop}
\begin{Shaded}
\begin{Highlighting}[]
\CommentTok{\# Create a cleaner dataset by dropping unnecessary columns}
\NormalTok{clean\_shipments }\OperatorTok{=}\NormalTok{ (}
\NormalTok{    shipments\_with\_lateness}
\NormalTok{    .drop(columns}\OperatorTok{=}\NormalTok{[}\StringTok{\textquotesingle{}quantity\textquotesingle{}}\NormalTok{])  }\CommentTok{\# Drop quantity column (not needed for our analysis)}
\NormalTok{    .dropna(subset}\OperatorTok{=}\NormalTok{[}\StringTok{\textquotesingle{}plannedShipDate\textquotesingle{}}\NormalTok{, }\StringTok{\textquotesingle{}actualShipDate\textquotesingle{}}\NormalTok{])  }\CommentTok{\# Remove rows with missing dates}
\NormalTok{)}

\BuiltInTok{print}\NormalTok{(}\SpecialStringTok{f"Cleaned dataset: }\SpecialCharTok{\{}\BuiltInTok{len}\NormalTok{(clean\_shipments)}\SpecialCharTok{\}}\SpecialStringTok{ rows, }\SpecialCharTok{\{}\BuiltInTok{len}\NormalTok{(clean\_shipments.columns)}\SpecialCharTok{\}}\SpecialStringTok{ columns"}\NormalTok{)}
\BuiltInTok{print}\NormalTok{(}\StringTok{"Remaining columns:"}\NormalTok{, clean\_shipments.columns.tolist())}
\end{Highlighting}
\end{Shaded}

\begin{verbatim}
Cleaned dataset: 4000 rows, 7 columns
Remaining columns: ['shipID', 'plannedShipDate', 'actualShipDate', 'partID', 'is_late', 'days_late', 'lateStatement']
\end{verbatim}

\begin{tcolorbox}[enhanced jigsaw, leftrule=.75mm, bottomtitle=1mm, toprule=.15mm, toptitle=1mm, left=2mm, colbacktitle=quarto-callout-important-color!10!white, rightrule=.15mm, titlerule=0mm, colback=white, colframe=quarto-callout-important-color-frame, title=\textcolor{quarto-callout-important-color}{\faExclamation}\hspace{0.5em}{🤔 Discussion Questions: Drop Mental Model}, arc=.35mm, coltitle=black, opacitybacktitle=0.6, breakable, opacityback=0, bottomrule=.15mm]

\textbf{Question 1: Drop vs Filter Strategies} - What's the difference
between
\texttt{.drop(columns={[}\textquotesingle{}quantity\textquotesingle{}{]})}
and \texttt{.filter()} with a list of columns you want to keep? - When
would you choose to drop columns vs filter to keep specific columns?

\textbf{Question 2: Handling Missing Data} - What happens if you use
\texttt{.dropna()} without specifying \texttt{subset}? How is this
different from
\texttt{.dropna(subset={[}\textquotesingle{}plannedShipDate\textquotesingle{},\ \textquotesingle{}actualShipDate\textquotesingle{}{]})}?
- Why might you want to be selective about which columns to check for
missing values?

\end{tcolorbox}

\paragraph{Briefly Give Answers to the Discussion Questions In This
Section}\label{briefly-give-answers-to-the-discussion-questions-in-this-section-2}

\begin{enumerate}
\def\labelenumi{\arabic{enumi}.}
\tightlist
\item
\end{enumerate}

Question 1: Drop vs Filter Strategies

.drop(columns={[}`quantity'{]}) Removes only the specified columns Keeps
all other columns Use when you want to remove a few specific columns
from many .filter({[}`col1', `col2'{]}) Keeps only the columns you list
Removes all other columns Use when you want to keep only a specific
subset of columns

\begin{Shaded}
\begin{Highlighting}[]
\CommentTok{\# These are equivalent for column selection:}
\NormalTok{method1 }\OperatorTok{=}\NormalTok{ shipments\_df.drop(columns}\OperatorTok{=}\NormalTok{[}\StringTok{\textquotesingle{}quantity\textquotesingle{}}\NormalTok{])  }\CommentTok{\# Explicit removal}
\NormalTok{method2 }\OperatorTok{=}\NormalTok{ shipments\_df.}\BuiltInTok{filter}\NormalTok{([}\StringTok{\textquotesingle{}shipID\textquotesingle{}}\NormalTok{, }\StringTok{\textquotesingle{}plannedShipDate\textquotesingle{}}\NormalTok{, }\StringTok{\textquotesingle{}actualShipDate\textquotesingle{}}\NormalTok{, }\StringTok{\textquotesingle{}partID\textquotesingle{}}\NormalTok{])  }\CommentTok{\# Explicit inclusion}

\BuiltInTok{print}\NormalTok{(}\StringTok{"Same columns:"}\NormalTok{, }\BuiltInTok{set}\NormalTok{(method1.columns) }\OperatorTok{==} \BuiltInTok{set}\NormalTok{(method2.columns))}
\end{Highlighting}
\end{Shaded}

\begin{verbatim}
Same columns: True
\end{verbatim}

Answer: Use .drop() when you know exactly which few columns to remove.
Use .filter() when you want to keep a specific subset of columns,
especially when you have many columns and only need a few.

2.To avoid losing good data unnecessarily

Example: A row has:

plannedShipDate: 2023-10-15 ✅

actualShipDate: 2023-10-14 ✅

customerNotes: NaN (missing)

.dropna() without subset: ❌ Removes the entire row because
customerNotes is missing

.dropna(subset={[}`plannedShipDate'{]}): ✅ Keeps the row because the
important date columns have data

Why be selective:

Keep rows where critical columns have data

Don't care if optional columns are missing

Preserve more data for analysis

Avoid deleting useful information

\begin{Shaded}
\begin{Highlighting}[]
\CommentTok{\# Without subset {-} drops rows with ANY missing values}
\NormalTok{dropped\_any\_na }\OperatorTok{=}\NormalTok{ shipments\_df.dropna()}
\BuiltInTok{print}\NormalTok{(}\SpecialStringTok{f"Rows after dropping any NA: }\SpecialCharTok{\{}\BuiltInTok{len}\NormalTok{(dropped\_any\_na)}\SpecialCharTok{\}}\SpecialStringTok{"}\NormalTok{)}

\CommentTok{\# With subset {-} only checks specific columns}
\NormalTok{dropped\_selected\_na }\OperatorTok{=}\NormalTok{ shipments\_df.dropna(subset}\OperatorTok{=}\NormalTok{[}\StringTok{\textquotesingle{}plannedShipDate\textquotesingle{}}\NormalTok{, }\StringTok{\textquotesingle{}actualShipDate\textquotesingle{}}\NormalTok{])}
\BuiltInTok{print}\NormalTok{(}\SpecialStringTok{f"Rows after dropping date NAs: }\SpecialCharTok{\{}\BuiltInTok{len}\NormalTok{(dropped\_selected\_na)}\SpecialCharTok{\}}\SpecialStringTok{"}\NormalTok{)}
\end{Highlighting}
\end{Shaded}

\begin{verbatim}
Rows after dropping any NA: 4000
Rows after dropping date NAs: 4000
\end{verbatim}

Using .dropna() without subset can remove valuable data unnecessarily.
Being selective with subset ensures we only remove rows where critical
fields are missing, preserving other useful data.

\subsubsection{4. Sort: Arranging Data}\label{sort-arranging-data}

\textbf{Mental Model:} Order data by values or indices.

Let's sort by lateness to see the worst offenders:

\phantomsection\label{mental-model-4-sort}
\begin{Shaded}
\begin{Highlighting}[]
\CommentTok{\# Sort by days late (worst first)}
\NormalTok{sorted\_by\_lateness }\OperatorTok{=}\NormalTok{ (}
\NormalTok{    clean\_shipments}
\NormalTok{    .sort\_values(}\StringTok{\textquotesingle{}days\_late\textquotesingle{}}\NormalTok{, ascending}\OperatorTok{=}\VariableTok{False}\NormalTok{)  }\CommentTok{\# Sort by days\_late, highest first}
\NormalTok{    .reset\_index(drop}\OperatorTok{=}\VariableTok{True}\NormalTok{)  }\CommentTok{\# Reset index to be sequential}
\NormalTok{)}

\BuiltInTok{print}\NormalTok{(}\StringTok{"Shipments sorted by lateness (worst first):"}\NormalTok{)}
\BuiltInTok{print}\NormalTok{(sorted\_by\_lateness[[}\StringTok{\textquotesingle{}shipID\textquotesingle{}}\NormalTok{, }\StringTok{\textquotesingle{}partID\textquotesingle{}}\NormalTok{, }\StringTok{\textquotesingle{}days\_late\textquotesingle{}}\NormalTok{, }\StringTok{\textquotesingle{}is\_late\textquotesingle{}}\NormalTok{]].head(}\DecValTok{10}\NormalTok{))}
\end{Highlighting}
\end{Shaded}

\begin{verbatim}
Shipments sorted by lateness (worst first):
   shipID       partID  days_late  is_late
0   10956  partb6208b5         21     True
1   10956  part04ef2f7         21     True
2   10956  part4875f85         21     True
3   10956  partb722d53         21     True
4   10956  partc979912         21     True
5   10956  parta27d449         21     True
6   10956  partc653823         21     True
7   10956  part82e69e9         21     True
8   10956  partf23fd1e         21     True
9   10956  part825873c         21     True
\end{verbatim}

\begin{tcolorbox}[enhanced jigsaw, leftrule=.75mm, bottomtitle=1mm, toprule=.15mm, toptitle=1mm, left=2mm, colbacktitle=quarto-callout-important-color!10!white, rightrule=.15mm, titlerule=0mm, colback=white, colframe=quarto-callout-important-color-frame, title=\textcolor{quarto-callout-important-color}{\faExclamation}\hspace{0.5em}{🤔 Discussion Questions: Sort Mental Model}, arc=.35mm, coltitle=black, opacitybacktitle=0.6, breakable, opacityback=0, bottomrule=.15mm]

\textbf{Question 1: Sorting Strategies} - What's the difference between
\texttt{ascending=False} and \texttt{ascending=True} in sorting? - How
would you sort by multiple columns (e.g., first by \texttt{is\_late},
then by \texttt{days\_late})?

\textbf{Question 2: Index Management} - Why do we use
\texttt{.reset\_index(drop=True)} after sorting? - What happens to the
original index when you sort? Why might this be problematic?

\end{tcolorbox}

\paragraph{Briefly Give Answers to the Discussion Questions In This
Section}\label{briefly-give-answers-to-the-discussion-questions-in-this-section-3}

\begin{enumerate}
\def\labelenumi{\arabic{enumi}.}
\tightlist
\item
  ascending=True (default):
\end{enumerate}

Sorts from lowest to highest

A-Z, 0-9, earliest to latest dates

Example: {[}1, 3, 5, 7{]}

ascending=False:

Sorts from highest to lowest

Z-A, 9-0, latest to earliest dates

Example: {[}7, 5, 3, 1{]}

\begin{Shaded}
\begin{Highlighting}[]
\CommentTok{\# Worst lateness first (highest days\_late at top)}
\NormalTok{shipments\_with\_lateness.sort\_values(}\StringTok{\textquotesingle{}days\_late\textquotesingle{}}\NormalTok{, ascending}\OperatorTok{=}\VariableTok{False}\NormalTok{)}

\CommentTok{\# Best lateness first (lowest days\_late at top)  }
\NormalTok{shipments\_with\_lateness.sort\_values(}\StringTok{\textquotesingle{}days\_late\textquotesingle{}}\NormalTok{, ascending}\OperatorTok{=}\VariableTok{True}\NormalTok{)}
\end{Highlighting}
\end{Shaded}

\begin{longtable}[]{@{}lllllllll@{}}
\toprule\noalign{}
& shipID & plannedShipDate & actualShipDate & partID & quantity &
is\_late & days\_late & lateStatement \\
\midrule\noalign{}
\endhead
\bottomrule\noalign{}
\endlastfoot
0 & 10001 & 2013-11-06 & 2013-10-04 & part92b16c5 & 6 & False & -33 &
Shipment is on Time \\
289 & 10059 & 2013-11-05 & 2013-10-11 & part6667ad & 4 & False & -25 &
Shipment is on Time \\
295 & 10059 & 2013-11-05 & 2013-10-11 & part0987f63 & 2 & False & -25 &
Shipment is on Time \\
296 & 10059 & 2013-11-05 & 2013-10-11 & partf194391 & 3 & False & -25 &
Shipment is on Time \\
293 & 10059 & 2013-11-05 & 2013-10-11 & partda5110b & 1 & False & -25 &
Shipment is on Time \\
... & ... & ... & ... & ... & ... & ... & ... & ... \\
3881 & 10956 & 2013-09-24 & 2013-10-15 & parta27d449 & 1 & True & 21 &
Darn Shipment is Late \\
3882 & 10956 & 2013-09-24 & 2013-10-15 & partc63f9bc & 1 & True & 21 &
Darn Shipment is Late \\
3883 & 10956 & 2013-09-24 & 2013-10-15 & part04ef2f7 & 1 & True & 21 &
Darn Shipment is Late \\
3884 & 10956 & 2013-09-24 & 2013-10-15 & part4875f85 & 1 & True & 21 &
Darn Shipment is Late \\
3894 & 10956 & 2013-09-24 & 2013-10-15 & part795d1a4 & 3 & True & 21 &
Darn Shipment is Late \\
\end{longtable}

1 b.

\begin{Shaded}
\begin{Highlighting}[]
\CommentTok{\# Sort by multiple columns}
\NormalTok{sorted\_df }\OperatorTok{=}\NormalTok{ shipments\_with\_lateness.sort\_values(}
\NormalTok{    [}\StringTok{\textquotesingle{}is\_late\textquotesingle{}}\NormalTok{, }\StringTok{\textquotesingle{}days\_late\textquotesingle{}}\NormalTok{], }
\NormalTok{    ascending}\OperatorTok{=}\NormalTok{[}\VariableTok{False}\NormalTok{, }\VariableTok{False}\NormalTok{]  }\CommentTok{\# is\_late: False, days\_late: False}
\NormalTok{)}
\BuiltInTok{print}\NormalTok{(sorted\_df[[}\StringTok{\textquotesingle{}shipID\textquotesingle{}}\NormalTok{, }\StringTok{\textquotesingle{}is\_late\textquotesingle{}}\NormalTok{, }\StringTok{\textquotesingle{}days\_late\textquotesingle{}}\NormalTok{]].head(}\DecValTok{10}\NormalTok{))}
\end{Highlighting}
\end{Shaded}

\begin{verbatim}
      shipID  is_late  days_late
1144   10217     True         21
3879   10956     True         21
3880   10956     True         21
3881   10956     True         21
3882   10956     True         21
3883   10956     True         21
3884   10956     True         21
3885   10956     True         21
3886   10956     True         21
3887   10956     True         21
\end{verbatim}

What this does:

First sorts by is\_late (False → True, so True appears first)

Then sorts by days\_late within each is\_late group (highest days\_late
first)

Result order:

Late shipments (is\_late=True) appear first

Within late shipments, most days late appear first

Then on-time shipments (is\_late=False)

Within on-time shipments, most days late appear first

2 a.

.reset\_index(drop=True) after sorting:

What it does:

Resets the row index to be sequential (0, 1, 2, 3\ldots)

drop=True removes the old index (doesn't keep it as a column)

Why we need it: After sorting, the original index gets messed up:

Index might be: {[}45, 12, 8, 33, 1\ldots{]} (random order)

After reset: {[}0, 1, 2, 3, 4\ldots{]} (clean order)

\begin{Shaded}
\begin{Highlighting}[]
\CommentTok{\# Without reset\_index {-} messy index}
\NormalTok{sorted\_messy }\OperatorTok{=}\NormalTok{ shipments\_with\_lateness.sort\_values(}\StringTok{\textquotesingle{}days\_late\textquotesingle{}}\NormalTok{, ascending}\OperatorTok{=}\VariableTok{False}\NormalTok{)}
\BuiltInTok{print}\NormalTok{(sorted\_messy.index)  }\CommentTok{\# [776, 777, 778, 779, 780...]}

\CommentTok{\# With reset\_index {-} clean index  }
\NormalTok{sorted\_clean }\OperatorTok{=}\NormalTok{ shipments\_with\_lateness.sort\_values(}\StringTok{\textquotesingle{}days\_late\textquotesingle{}}\NormalTok{, ascending}\OperatorTok{=}\VariableTok{False}\NormalTok{).reset\_index(drop}\OperatorTok{=}\VariableTok{True}\NormalTok{)}
\BuiltInTok{print}\NormalTok{(sorted\_clean.index)  }\CommentTok{\# [0, 1, 2, 3, 4...]}
\end{Highlighting}
\end{Shaded}

\begin{verbatim}
Index([3888, 3883, 3884, 3885, 3886, 3881, 3887, 3897, 3898, 3899,
       ...
        299,  298,  293,  304,  324,  323,  320,  321,  322,    0],
      dtype='int64', length=4000)
RangeIndex(start=0, stop=4000, step=1)
\end{verbatim}

Why this matters:

Clean indexes are easier to work with

Prevents confusion when referencing rows

Some operations expect sequential indexes

2b.

What happens to original index when sorting:

The index stays with the same rows but gets shuffled

Example: Row that was at index 5 might move to index 150 after sorting

Why this is problematic:

Confusing row references:

::: \{.cell execution\_count=18\}
\texttt{\{.python\ .cell-code\}\ \ \#\ After\ sorting,\ index\ 0\ might\ not\ be\ the\ first\ original\ row\ \ sorted\_df.iloc{[}0{]}\ \#\ This\ might\ have\ been\ row\ 45\ in\ original\ data}

::: \{.cell-output .cell-output-display execution\_count=73\}
\texttt{shipID\ \ \ \ \ \ \ \ \ \ \ \ \ \ \ \ \ \ \ \ \ \ \ \ \ \ \ \ \ 10217\ \ plannedShipDate\ \ \ \ \ \ 2013-09-23\ 00:00:00\ \ actualShipDate\ \ \ \ \ \ \ 2013-10-14\ 00:00:00\ \ partID\ \ \ \ \ \ \ \ \ \ \ \ \ \ \ \ \ \ \ \ \ \ \ part2081be9\ \ quantity\ \ \ \ \ \ \ \ \ \ \ \ \ \ \ \ \ \ \ \ \ \ \ \ \ \ \ \ \ \ \ 5\ \ is\_late\ \ \ \ \ \ \ \ \ \ \ \ \ \ \ \ \ \ \ \ \ \ \ \ \ \ \ \ \ True\ \ days\_late\ \ \ \ \ \ \ \ \ \ \ \ \ \ \ \ \ \ \ \ \ \ \ \ \ \ \ \ \ 21\ \ lateStatement\ \ \ \ \ \ Darn\ Shipment\ is\ Late\ \ Name:\ 1144,\ dtype:\ object}
::: :::

Broken relationships:

If you had external references to row indexes, they're now wrong

Messy display:

Index shows random numbers: {[}45, 12, 8, 33, 1\ldots{]}

Hard to see the actual sorted order

Some functions expect sequential indexes:

Plotting, some pandas operations work better with clean 0,1,2,3\ldots{}
indexes

\subsubsection{5. Aggregate: Summarizing
Data}\label{aggregate-summarizing-data}

\textbf{Mental Model:} Calculate summary statistics across groups or the
entire dataset.

Let's calculate overall service level metrics:

\phantomsection\label{mental-model-5-aggregate}
\begin{Shaded}
\begin{Highlighting}[]
\CommentTok{\# Calculate overall service level metrics}
\NormalTok{service\_metrics }\OperatorTok{=}\NormalTok{ (}
\NormalTok{    clean\_shipments}
\NormalTok{    .agg(\{}
        \StringTok{\textquotesingle{}is\_late\textquotesingle{}}\NormalTok{: [}\StringTok{\textquotesingle{}count\textquotesingle{}}\NormalTok{, }\StringTok{\textquotesingle{}sum\textquotesingle{}}\NormalTok{, }\StringTok{\textquotesingle{}mean\textquotesingle{}}\NormalTok{],  }\CommentTok{\# Count total, count late, calculate percentage}
        \StringTok{\textquotesingle{}days\_late\textquotesingle{}}\NormalTok{: [}\StringTok{\textquotesingle{}mean\textquotesingle{}}\NormalTok{, }\StringTok{\textquotesingle{}max\textquotesingle{}}\NormalTok{]  }\CommentTok{\# Average and maximum days late}
\NormalTok{    \})}
\NormalTok{    .}\BuiltInTok{round}\NormalTok{(}\DecValTok{3}\NormalTok{)}
\NormalTok{)}

\BuiltInTok{print}\NormalTok{(}\StringTok{"Overall Service Level Metrics:"}\NormalTok{)}
\BuiltInTok{print}\NormalTok{(service\_metrics)}

\CommentTok{\# Calculate percentage on{-}time directly from the data}
\NormalTok{on\_time\_rate }\OperatorTok{=}\NormalTok{ (}\DecValTok{1} \OperatorTok{{-}}\NormalTok{ clean\_shipments[}\StringTok{\textquotesingle{}is\_late\textquotesingle{}}\NormalTok{].mean()) }\OperatorTok{*} \DecValTok{100}
\BuiltInTok{print}\NormalTok{(}\SpecialStringTok{f"}\CharTok{\textbackslash{}n}\SpecialStringTok{On{-}time delivery rate: }\SpecialCharTok{\{}\NormalTok{on\_time\_rate}\SpecialCharTok{:.1f\}}\SpecialStringTok{\%"}\NormalTok{)}
\end{Highlighting}
\end{Shaded}

\begin{verbatim}
Overall Service Level Metrics:
        is_late  days_late
count  4000.000        NaN
sum     456.000        NaN
mean      0.114     -0.974
max         NaN     21.000

On-time delivery rate: 88.6%
\end{verbatim}

\begin{tcolorbox}[enhanced jigsaw, leftrule=.75mm, bottomtitle=1mm, toprule=.15mm, toptitle=1mm, left=2mm, colbacktitle=quarto-callout-important-color!10!white, rightrule=.15mm, titlerule=0mm, colback=white, colframe=quarto-callout-important-color-frame, title=\textcolor{quarto-callout-important-color}{\faExclamation}\hspace{0.5em}{🤔 Discussion Questions: Aggregate Mental Model}, arc=.35mm, coltitle=black, opacitybacktitle=0.6, breakable, opacityback=0, bottomrule=.15mm]

\textbf{Question 1: Boolean Aggregation} - Why does \texttt{sum()} work
on boolean values? What does it count?

\end{tcolorbox}

\paragraph{Briefly Give Answers to the Discussion Questions In This
Section}\label{briefly-give-answers-to-the-discussion-questions-in-this-section-4}

Why sum() works on boolean values:

In Python/pandas, True = 1 and False = 0

What sum() counts:

It counts how many True values (1's) are in the boolean series

Each True = 1, each False = 0

\begin{Shaded}
\begin{Highlighting}[]
\CommentTok{\# First create clean\_shipments (from earlier cells)}
\NormalTok{clean\_shipments }\OperatorTok{=}\NormalTok{ (}
\NormalTok{    shipments\_with\_lateness}
\NormalTok{    .drop(columns}\OperatorTok{=}\NormalTok{[}\StringTok{\textquotesingle{}quantity\textquotesingle{}}\NormalTok{])  }\CommentTok{\# Drop quantity column}
\NormalTok{    .dropna(subset}\OperatorTok{=}\NormalTok{[}\StringTok{\textquotesingle{}plannedShipDate\textquotesingle{}}\NormalTok{, }\StringTok{\textquotesingle{}actualShipDate\textquotesingle{}}\NormalTok{])  }\CommentTok{\# Remove rows with missing dates}
\NormalTok{)}

\CommentTok{\# Then calculate service metrics}
\NormalTok{service\_metrics }\OperatorTok{=}\NormalTok{ (}
\NormalTok{    clean\_shipments}
\NormalTok{    .agg(\{}
        \StringTok{\textquotesingle{}is\_late\textquotesingle{}}\NormalTok{: [}\StringTok{\textquotesingle{}count\textquotesingle{}}\NormalTok{, }\StringTok{\textquotesingle{}sum\textquotesingle{}}\NormalTok{, }\StringTok{\textquotesingle{}mean\textquotesingle{}}\NormalTok{],  }\CommentTok{\# Count total, count late, calculate percentage}
        \StringTok{\textquotesingle{}days\_late\textquotesingle{}}\NormalTok{: [}\StringTok{\textquotesingle{}mean\textquotesingle{}}\NormalTok{, }\StringTok{\textquotesingle{}max\textquotesingle{}}\NormalTok{]  }\CommentTok{\# Average and maximum days late}
\NormalTok{    \})}
\NormalTok{    .}\BuiltInTok{round}\NormalTok{(}\DecValTok{3}\NormalTok{)}
\NormalTok{)}

\BuiltInTok{print}\NormalTok{(}\StringTok{"Overall Service Level Metrics:"}\NormalTok{)}
\BuiltInTok{print}\NormalTok{(service\_metrics)}

\CommentTok{\# Calculate percentage on{-}time directly from the data}
\NormalTok{on\_time\_rate }\OperatorTok{=}\NormalTok{ (}\DecValTok{1} \OperatorTok{{-}}\NormalTok{ clean\_shipments[}\StringTok{\textquotesingle{}is\_late\textquotesingle{}}\NormalTok{].mean()) }\OperatorTok{*} \DecValTok{100}
\BuiltInTok{print}\NormalTok{(}\SpecialStringTok{f"}\CharTok{\textbackslash{}n}\SpecialStringTok{On{-}time delivery rate: }\SpecialCharTok{\{}\NormalTok{on\_time\_rate}\SpecialCharTok{:.1f\}}\SpecialStringTok{\%"}\NormalTok{)}
\end{Highlighting}
\end{Shaded}

\begin{verbatim}
Overall Service Level Metrics:
        is_late  days_late
count  4000.000        NaN
sum     456.000        NaN
mean      0.114     -0.974
max         NaN     21.000

On-time delivery rate: 88.6%
\end{verbatim}

In our service metrics:

\begin{itemize}
\tightlist
\item
  \texttt{\textquotesingle{}is\_late\textquotesingle{}:\ {[}\textquotesingle{}count\textquotesingle{},\ \textquotesingle{}sum\textquotesingle{},\ \textquotesingle{}mean\textquotesingle{}{]}}
\item
  count: Total number of shipments (4000)
\item
  sum: Number of late shipments (456 True values)
\item
  mean: Percentage of late shipments (456/4000 = 0.114)
\end{itemize}

So sum() on boolean values gives you the count of True values.

\subsubsection{6. Merge: Combining
Information}\label{merge-combining-information}

\textbf{Mental Model:} Join data from multiple sources to create richer
datasets.

Now let's analyze service levels by product category. First, we need to
merge our data:

\phantomsection\label{mental-model-6-merge-prep}
\begin{Shaded}
\begin{Highlighting}[]
\CommentTok{\# Merge shipment data with product line data}
\NormalTok{shipments\_with\_category }\OperatorTok{=}\NormalTok{ (}
\NormalTok{    clean\_shipments}
\NormalTok{    .merge(product\_line\_df, on}\OperatorTok{=}\StringTok{\textquotesingle{}partID\textquotesingle{}}\NormalTok{, how}\OperatorTok{=}\StringTok{\textquotesingle{}left\textquotesingle{}}\NormalTok{)}
\NormalTok{    .assign(}
\NormalTok{        category\_late}\OperatorTok{=}\KeywordTok{lambda}\NormalTok{ df: df[}\StringTok{\textquotesingle{}is\_late\textquotesingle{}}\NormalTok{] }\OperatorTok{\&}\NormalTok{ df[}\StringTok{\textquotesingle{}prodCategory\textquotesingle{}}\NormalTok{].notna()}
\NormalTok{    )}
\NormalTok{)}

\BuiltInTok{print}\NormalTok{(}\StringTok{"}\CharTok{\textbackslash{}n}\StringTok{Product categories available:"}\NormalTok{)}
\BuiltInTok{print}\NormalTok{(shipments\_with\_category[}\StringTok{\textquotesingle{}prodCategory\textquotesingle{}}\NormalTok{].value\_counts())}
\end{Highlighting}
\end{Shaded}

\begin{verbatim}

Product categories available:
prodCategory
Marketables    1850
Machines        846
SpareParts      767
Liquids         537
Name: count, dtype: int64
\end{verbatim}

\begin{tcolorbox}[enhanced jigsaw, leftrule=.75mm, bottomtitle=1mm, toprule=.15mm, toptitle=1mm, left=2mm, colbacktitle=quarto-callout-important-color!10!white, rightrule=.15mm, titlerule=0mm, colback=white, colframe=quarto-callout-important-color-frame, title=\textcolor{quarto-callout-important-color}{\faExclamation}\hspace{0.5em}{🤔 Discussion Questions: Merge Mental Model}, arc=.35mm, coltitle=black, opacitybacktitle=0.6, breakable, opacityback=0, bottomrule=.15mm]

\textbf{Question 1: Join Types and Data Loss} - Why does your professor
think we should use
\texttt{how=\textquotesingle{}left\textquotesingle{}} in most cases? -
How can you check if any shipments were lost during the merge?

\textbf{Question 2: Key Column Matching} - What happens if there are
duplicate \texttt{partID} values in the \texttt{product\_line\_df}?

\end{tcolorbox}

\paragraph{Briefly Give Answers to the Discussion Questions In This
Section}\label{briefly-give-answers-to-the-discussion-questions-in-this-section-5}

Question 1: Why use how=`left'

Prevents data loss - keeps all shipments even if category data is
missing

Shows data gaps clearly with NaN values

Preserves the primary dataset

Question 1b: Check for lost shipments

Question 2: Duplicate partIDs

Creates duplicate rows in results

Inflates counts and statistics

Check with: product\_line\_df.duplicated(`partID').sum()

\subsubsection{7. Split-Apply-Combine: Group
Analysis}\label{split-apply-combine-group-analysis}

\textbf{Mental Model:} Group data and apply functions within each group.

Now let's analyze service levels by category:

\phantomsection\label{mental-model-7-groupby}
\begin{Shaded}
\begin{Highlighting}[]
\CommentTok{\# Analyze service levels by product category}
\NormalTok{service\_by\_category }\OperatorTok{=}\NormalTok{ (}
\NormalTok{    shipments\_with\_category}
\NormalTok{    .groupby(}\StringTok{\textquotesingle{}prodCategory\textquotesingle{}}\NormalTok{)  }\CommentTok{\# Split by product category}
\NormalTok{    .agg(\{}
        \StringTok{\textquotesingle{}is\_late\textquotesingle{}}\NormalTok{: [}\StringTok{\textquotesingle{}any\textquotesingle{}}\NormalTok{, }\StringTok{\textquotesingle{}count\textquotesingle{}}\NormalTok{, }\StringTok{\textquotesingle{}sum\textquotesingle{}}\NormalTok{, }\StringTok{\textquotesingle{}mean\textquotesingle{}}\NormalTok{],  }\CommentTok{\# Count, late count, percentage late}
        \StringTok{\textquotesingle{}days\_late\textquotesingle{}}\NormalTok{: [}\StringTok{\textquotesingle{}mean\textquotesingle{}}\NormalTok{, }\StringTok{\textquotesingle{}max\textquotesingle{}}\NormalTok{]  }\CommentTok{\# Average and max days late}
\NormalTok{    \})}
\NormalTok{    .}\BuiltInTok{round}\NormalTok{(}\DecValTok{3}\NormalTok{)}
\NormalTok{)}

\BuiltInTok{print}\NormalTok{(}\StringTok{"Service Level by Product Category:"}\NormalTok{)}
\BuiltInTok{print}\NormalTok{(service\_by\_category)}
\end{Highlighting}
\end{Shaded}

\begin{verbatim}
Service Level by Product Category:
             is_late                   days_late    
                 any count  sum   mean      mean max
prodCategory                                        
Liquids         True   537   22  0.041    -0.950  19
Machines        True   846  152  0.180    -1.336  21
Marketables     True  1850  145  0.078    -0.804  21
SpareParts      True   767  137  0.179    -1.003  21
\end{verbatim}

\begin{tcolorbox}[enhanced jigsaw, leftrule=.75mm, bottomtitle=1mm, toprule=.15mm, toptitle=1mm, left=2mm, colbacktitle=quarto-callout-important-color!10!white, rightrule=.15mm, titlerule=0mm, colback=white, colframe=quarto-callout-important-color-frame, title=\textcolor{quarto-callout-important-color}{\faExclamation}\hspace{0.5em}{🤔 Discussion Questions: Split-Apply-Combine Mental Model}, arc=.35mm, coltitle=black, opacitybacktitle=0.6, breakable, opacityback=0, bottomrule=.15mm]

\textbf{Question 1: GroupBy Mechanics} - What does
\texttt{.groupby(\textquotesingle{}prodCategory\textquotesingle{})}
actually do? How does it ``split'' the data? - Why do we need to use
\texttt{.agg()} after grouping? What happens if you don't?

\textbf{Question 2: Multi-Level Grouping} - Explore grouping by
\texttt{{[}\textquotesingle{}shipID\textquotesingle{},\ \textquotesingle{}prodCategory\textquotesingle{}{]}}?
What question does this answer versus grouping by
\texttt{\textquotesingle{}prodCategory\textquotesingle{}} alone? (HINT:
There may be many rows with identical shipID's due to a particular order
having multiple partID's.)

\end{tcolorbox}

\paragraph{Briefly Give Answers to the Discussion Questions In This
Section}\label{briefly-give-answers-to-the-discussion-questions-in-this-section-6}

1 a.

What .groupby(`prodCategory') does:

It ``splits'' the data into separate groups based on unique category
values: Before grouping:

text shipID partID prodCategory is\_late days\_late 10001 part92b16c5
Marketables False -33 10002 part66983b Machines False -11\\
10003 part8e36f25 Marketables False -18 10004 part30f5de0 Liquids False
-6 10005 part9d64d35 SpareParts False -6 After .groupby(`prodCategory'):

Group 1: Marketables

text 10001 part92b16c5 Marketables False -33 10003 part8e36f25
Marketables False -18 \ldots{} Group 2: Machines

text 10002 part66983b Machines False -11 \ldots{} Group 3: Liquids

text 10004 part30f5de0 Liquids False -6 \ldots{} Group 4: SpareParts

text 10005 part9d64d35 SpareParts False -6 \ldots{} How it works:
Creates separate ``mini-dataframes'' for each category

Each group contains only rows from that category

When you apply .agg(), it calculates statistics within each group
separately

The ``split-apply-combine'' pattern:

Split data into groups

Apply functions to each group

Combine results back together

1 b.

Why we need .agg() after grouping:

Without .agg(): The grouped data is incomplete and unusable - it's just
split into groups but no calculations are performed.

::: \{.cell execution\_count=23\}
\texttt{\{.python\ .cell-code\}\ \ \#\ This\ creates\ a\ DataFrameGroupBy\ object,\ not\ useful\ results\ \ grouped=\ shipments\_with\_category.groupby(\textquotesingle{}prodCategory\textquotesingle{})\ \ print(type(grouped))\ \#\ \textless{}class\ \textquotesingle{}pandas.core.groupby.generic.DataFrameGroupBy\textquotesingle{}\textgreater{}\ \ print(grouped)\ \ \ \ \ \ \ \#\ Can\textquotesingle{}t\ see\ the\ actual\ data}

::: \{.cell-output .cell-output-stdout\}
\texttt{\textless{}class\ \textquotesingle{}pandas.core.groupby.generic.DataFrameGroupBy\textquotesingle{}\textgreater{}\ \ \textless{}pandas.core.groupby.generic.DataFrameGroupBy\ object\ at\ 0x000002588980F320\textgreater{}}
::: :::

What .agg() does: It applies calculations to each group and combines the
results:

::: \{.cell execution\_count=24\} ``` \{.python .cell-code\} \# Without
.agg() - useless grouped\_data=
shipments\_with\_category.groupby(`prodCategory') \#Can't see or use
this meaningfully

\#With .agg() - useful results results=
shipments\_with\_category.groupby(`prodCategory').agg(\{ `is\_late':
`mean', `days\_late': `max' \}) \#Gets: average lateness and max days
late FOR EACH CATEGORY ``` :::

What happens if you don't use .agg(): You get a GroupBy object, not a
DataFrame

No calculations are performed

You can't see or analyze the grouped results

The data remains ``split'' but not summarized

Bottom line: .agg() is what transforms the split groups into meaningful
summary statistics you can actually use and analyze.

\begin{enumerate}
\def\labelenumi{\arabic{enumi}.}
\setcounter{enumi}{1}
\tightlist
\item
\end{enumerate}

Grouping by {[}`shipID', `prodCategory'{]} vs `prodCategory' alone:

Grouping by `prodCategory' alone: Answers: ``What is the service level
for each product category?''

Looks at individual line items

Doesn't care which shipment they belong to

Example: ``4\% of Liquids line items are late''

::: \{.cell execution\_count=25\}
\texttt{\{.python\ .cell-code\}\ \ service\_by\_category\ =\ (\ \ \ \ \ shipments\_with\_category\ \ \ \ \ .groupby(\textquotesingle{}prodCategory\textquotesingle{})\ \ \ \ \ .agg(\{\textquotesingle{}is\_late\textquotesingle{}:\ \textquotesingle{}mean\textquotesingle{}\})\ \ )}
:::

Grouping by {[}`shipID', `prodCategory'{]}: Answers: ``For shipments
containing multiple categories, how does lateness vary between
categories within the same shipment?''

Looks at shipment-category combinations

Recognizes that one shipment can have items from multiple categories

Example: ``In shipment 10192, were the Machines items late but the
Liquids on time?''

::: \{.cell execution\_count=26\} ``` \{.python .cell-code\}
shipment\_category\_analysis= ( shipments\_with\_category
.groupby({[}`shipID', `prodCategory'{]}) .agg(\{`is\_late': `any'\}) \#
Was ANY item in this shipment-category combo late? )

\begin{verbatim}
:::


Key difference:
prodCategory alone: "How late is each category overall?"

[shipID, prodCategory]: "Within multi-category shipments, which categories tend to be the problem?"

This helps answer: "When an order has both Machines and Liquids, is one category consistently causing the delay?"

## Answering A Business Question

**Mental Model:** Combine multiple data manipulation techniques to answer complex business questions.

Let's create a comprehensive analysis by combining shipment-level data with category information:

::: {#mental-model-7-comprehensive .cell execution_count=27}
``` {.python .cell-code}
# Create a comprehensive analysis dataset
comprehensive_analysis = (
   shipments_with_category
   .groupby(['shipID', 'prodCategory'])  # Group by shipment and category
   .agg({
       'is_late': 'any',  # True if any item in this shipment/category is late
       'days_late': 'max'  # Maximum days late for this shipment/category
   })
   .reset_index()
   .assign(
       has_multiple_categories=lambda df: df.groupby('shipID')['prodCategory'].transform('nunique') > 1
   )
)

print("Comprehensive analysis - shipments with multiple categories:")
multi_category_shipments = comprehensive_analysis[comprehensive_analysis['has_multiple_categories']]
print(f"Shipments with multiple categories: {multi_category_shipments['shipID'].nunique()}")
print(f"Total unique shipments: {comprehensive_analysis['shipID'].nunique()}")
print(f"Percentage with multiple categories: {multi_category_shipments['shipID'].nunique() / comprehensive_analysis['shipID'].nunique() * 100:.1f}%")
\end{verbatim}

\begin{verbatim}
Comprehensive analysis - shipments with multiple categories:
Shipments with multiple categories: 232
Total unique shipments: 997
Percentage with multiple categories: 23.3%
\end{verbatim}

:::

\begin{tcolorbox}[enhanced jigsaw, leftrule=.75mm, bottomtitle=1mm, toprule=.15mm, toptitle=1mm, left=2mm, colbacktitle=quarto-callout-important-color!10!white, rightrule=.15mm, titlerule=0mm, colback=white, colframe=quarto-callout-important-color-frame, title=\textcolor{quarto-callout-important-color}{\faExclamation}\hspace{0.5em}{🤔 Discussion Questions: Answering A Business Question}, arc=.35mm, coltitle=black, opacitybacktitle=0.6, breakable, opacityback=0, bottomrule=.15mm]

\textbf{Question 1: Business Question Analysis} - What business question
does this comprehensive analysis answer? - How does grouping by
\texttt{{[}\textquotesingle{}shipID\textquotesingle{},\ \textquotesingle{}prodCategory\textquotesingle{}{]}}
differ from grouping by just
\texttt{\textquotesingle{}prodCategory\textquotesingle{}}? - What
insights can ZappTech's management gain from knowing the percentage of
multi-category shipments?

\end{tcolorbox}

\paragraph{Briefly Give Answers to the Discussion Questions In This
Section}\label{briefly-give-answers-to-the-discussion-questions-in-this-section-7}

\begin{enumerate}
\def\labelenumi{\arabic{enumi}.}
\tightlist
\item
\end{enumerate}

``How common are multi-category orders, and how does lateness vary
within them?'' Specifically:

``What percentage of shipments contain products from multiple
categories?'' (23.3\%)

``When orders span multiple categories, which categories within the same
shipment are late vs on-time?''

``Do certain category combinations in the same order have different
service levels?''

Business insights for ZappTech management: Customer Behavior:

23.3\% of customers order across multiple product categories in the same
shipment

This reveals cross-selling patterns and customer purchasing habits

Operational Analysis:

Within multi-category shipments, management can see if specific
categories consistently cause delays

Example: ``When orders contain both Machines and Liquids, are the
Machines always late while Liquids are on time?''

Supply Chain Optimization:

Identifies which category combinations create the most operational
complexity

Helps optimize warehouse picking, packing, and shipping for
mixed-category orders

Service Level Management:

Provides granular insight into service levels at the shipment-category
level, not just overall category level

Helps pinpoint exactly where in the fulfillment process delays occur for
complex orders

\begin{enumerate}
\def\labelenumi{\arabic{enumi}.}
\setcounter{enumi}{1}
\tightlist
\item
\end{enumerate}

Grouping by {[}`shipID', `prodCategory'{]} vs `prodCategory':

Grouping by `prodCategory' alone: Answers: ``What is the overall service
level for each category?''

Perspective: Category-level analysis

Example: ``Machines have 18\% late rate across all shipments''

python \# Result: One row per category prodCategory is\_late\_mean
Machines 0.180 Liquids 0.041 Marketables 0.078 Grouping by {[}`shipID',
`prodCategory'{]}: Answers: ``Within each multi-category shipment, how
does service level vary between categories?''

Perspective: Shipment-level analysis across categories

Example: ``In shipment 10192, Machines were late but Liquids were on
time''

\section{Result: One row per shipment-category
combination}\label{result-one-row-per-shipment-category-combination}

shipID prodCategory is\_late\_any 10192 Machines True 10192 Liquids
False 10193 Marketables True\\
10193 SpareParts False Key difference: prodCategory: ``How is each
category performing overall?''

This reveals intra-shipment patterns - how different categories behave
when they're part of the same customer order versus when they're ordered
separately.

\begin{enumerate}
\def\labelenumi{\arabic{enumi}.}
\setcounter{enumi}{2}
\tightlist
\item
\end{enumerate}

Business insights from multi-category shipment percentage (23.3\%):

Customer Behavior \& Marketing: Cross-selling success: Nearly 1 in 4
customers buy across categories in single orders

Bundling opportunities: Identify which category pairs are frequently
purchased together

Marketing strategy: Target customers with complementary products from
different categories

Operations \& Logistics: Warehouse optimization: 23.3\% of shipments
require picking from multiple category zones

Staff training: Workers need to handle diverse product types in single
shipments

Packaging complexity: Mixed-category orders may need special
handling/packaging

Supply Chain Management: Inventory planning: Stock complementary
categories in proximity

Shipping efficiency: Multi-category orders might take longer to assemble

Quality control: Different categories may have different handling
requirements

Customer Experience: Order fulfillment: Multi-category orders risk
partial shipments if one category is unavailable

Delivery expectations: Customers expect all categories to arrive
together

Service level: Late items in one category can make the entire shipment
appear late

Strategic value: This 23.3\% represents ZappTech's most valuable and
complex orders - customers who see the company as a one-stop shop rather
than a single-category supplier.

\subsection{Student Analysis Section: Mastering Data
Manipulation}\label{student-analysis-section}

\textbf{Your Task:} Demonstrate your mastery of the seven mental models
through comprehensive discussion and analysis. The bulk of your grade
comes from thoughtfully answering the discussion questions for each
mental model. See below for more details.

\subsubsection{Core Challenge: Discussion Questions
Analysis}\label{core-challenge-discussion-questions-analysis}

\textbf{For each mental model, provide:} - Clear, concise answers to all
discussion questions - Code examples where appropriate to support your
explanations

\begin{tcolorbox}[enhanced jigsaw, leftrule=.75mm, bottomtitle=1mm, toprule=.15mm, toptitle=1mm, left=2mm, colbacktitle=quarto-callout-important-color!10!white, rightrule=.15mm, titlerule=0mm, colback=white, colframe=quarto-callout-important-color-frame, title=\textcolor{quarto-callout-important-color}{\faExclamation}\hspace{0.5em}{📊 Discussion Questions Requirements}, arc=.35mm, coltitle=black, opacitybacktitle=0.6, breakable, opacityback=0, bottomrule=.15mm]

\textbf{Complete all discussion question sections:} 1. \textbf{Assign
Mental Model:} Data types, date handling, and debugging 2.
\textbf{Subset Mental Model:} Filtering strategies and complex queries
3. \textbf{Drop Mental Model:} Data cleaning and quality management 4.
\textbf{Sort Mental Model:} Data organization and business logic 5.
\textbf{Aggregate Mental Model:} Summary statistics and business metrics
6. \textbf{Merge Mental Model:} Data integration and quality control 7.
\textbf{Split-Apply-Combine Mental Model:} Group analysis and advanced
operations 8. \textbf{Answering A Business Question:} Combining multiple
data manipulation techniques to answer a business question

\end{tcolorbox}

\subsubsection{Professional Visualizations (For 100\%
Grade)}\label{professional-visualizations-for-100-grade}

\textbf{Your Task:} Create a professional visualization that supports
your analysis and demonstrates your understanding of the data.

import matplotlib.pyplot as plt import seaborn as sns

\section{First, ensure we have all required data by recreating the
necessary
variables}\label{first-ensure-we-have-all-required-data-by-recreating-the-necessary-variables}

\section{Load data if not already
loaded}\label{load-data-if-not-already-loaded}

if `shipments\_df' not in globals(): shipments\_df = pd.read\_csv(
``https://raw.githubusercontent.com/flyaflya/persuasive/main/shipments.csv'',
parse\_dates={[}`plannedShipDate', `actualShipDate'{]} ).head(4000)

if `product\_line\_df' not in globals(): product\_line\_df =
pd.read\_csv(
``https://raw.githubusercontent.com/flyaflya/persuasive/main/productLine.csv''
)

\section{Create shipments\_with\_lateness if not
exists}\label{create-shipments_with_lateness-if-not-exists}

if `shipments\_with\_lateness' not in globals():
shipments\_with\_lateness = ( shipments\_df .assign( is\_late=lambda df:
df{[}`actualShipDate'{]} \textgreater{} df{[}`plannedShipDate'{]},
days\_late=lambda df: (df{[}`actualShipDate'{]} -
df{[}`plannedShipDate'{]}).dt.days ) )

\section{Create clean\_shipments if not
exists}\label{create-clean_shipments-if-not-exists}

if `clean\_shipments' not in globals(): clean\_shipments = (
shipments\_with\_lateness .drop(columns={[}`quantity'{]})
.dropna(subset={[}`plannedShipDate', `actualShipDate'{]}) )

\section{Create shipments\_with\_category if not
exists}\label{create-shipments_with_category-if-not-exists}

if `shipments\_with\_category' not in globals():
shipments\_with\_category = ( clean\_shipments .merge(product\_line\_df,
on=`partID', how=`left') .assign( category\_late=lambda df:
df{[}`is\_late'{]} \& df{[}`prodCategory'{]}.notna() ) )

\section{Create comprehensive\_analysis if not
exists}\label{create-comprehensive_analysis-if-not-exists}

if `comprehensive\_analysis' not in globals(): comprehensive\_analysis =
( shipments\_with\_category .groupby({[}`shipID', `prodCategory'{]})
.agg(\{ `is\_late': `any', `days\_late': `max' \}) .reset\_index()
.assign( has\_multiple\_categories=lambda df:
df.groupby(`shipID'){[}`prodCategory'{]}.transform(`nunique')
\textgreater{} 1 ) )

\section{Create multi\_category\_shipments if not
exists}\label{create-multi_category_shipments-if-not-exists}

if `multi\_category\_shipments' not in globals():
multi\_category\_shipments =
comprehensive\_analysis{[}comprehensive\_analysis{[}`has\_multiple\_categories'{]}{]}

\section{Calculate overall on-time rate if not
exists}\label{calculate-overall-on-time-rate-if-not-exists}

if `on\_time\_rate' not in globals(): on\_time\_rate = (1 -
clean\_shipments{[}`is\_late'{]}.mean()) * 100

\section{Calculate service level metrics by
category}\label{calculate-service-level-metrics-by-category}

category\_service = ( shipments\_with\_category .groupby(`prodCategory')
.agg(\{ `is\_late': {[}`count', `sum', `mean'{]}, `days\_late': `mean'
\}) .round(3) ) category\_service.columns = {[}`total\_shipments',
`late\_count', `late\_rate', `avg\_days\_late'{]}
category\_service{[}`on\_time\_rate'{]} = (1 -
category\_service{[}`late\_rate'{]}) * 100

\section{Create comprehensive professional visualization demonstrating
data
understanding}\label{create-comprehensive-professional-visualization-demonstrating-data-understanding}

fig = plt.figure(figsize=(16, 12)) gs = fig.add\_gridspec(2, 3,
height\_ratios={[}1, 1{]}, width\_ratios={[}1, 1, 1{]})

\section{Professional color palette}\label{professional-color-palette}

colors = {[}`\#2E8B57', `\#4682B4', `\#DAA520', `\#CD5C5C'{]} \# Green,
Blue, Gold, Red plt.style.use(`default')

\section{Plot 1: On-time delivery rate by category (Top
Left)}\label{plot-1-on-time-delivery-rate-by-category-top-left}

ax1 = fig.add\_subplot(gs{[}0, 0{]}) bars1 =
ax1.bar(category\_service.index,
category\_service{[}`on\_time\_rate'{]}, color=colors, alpha=0.8,
edgecolor=`black', linewidth=1) ax1.set\_title(`On-Time Delivery Rate by
Category', fontsize=12, fontweight=`bold', pad=15)
ax1.set\_ylabel(`On-Time Rate (\%)', fontweight=`bold')
ax1.set\_ylim(75, 100) ax1.grid(axis=`y', alpha=0.3, linestyle=`--')

\section{Add value labels}\label{add-value-labels}

for bar, value in zip(bars1, category\_service{[}`on\_time\_rate'{]}):
ax1.text(bar.get\_x() + bar.get\_width()/2, bar.get\_height() + 0.3,
f'\{value:.1f\}\%`, ha='center', va=`bottom', fontweight=`bold',
fontsize=10)

\section{Plot 2: Average days late by category (Top
Center)}\label{plot-2-average-days-late-by-category-top-center}

ax2 = fig.add\_subplot(gs{[}0, 1{]}) bars2 =
ax2.bar(category\_service.index,
category\_service{[}`avg\_days\_late'{]}, color=colors, alpha=0.8,
edgecolor=`black', linewidth=1) ax2.set\_title(`Average Days Late by
Category', fontsize=12, fontweight=`bold', pad=15)
ax2.set\_ylabel(`Average Days Late', fontweight=`bold')
ax2.grid(axis=`y', alpha=0.3, linestyle=`--')

\section{Add value labels}\label{add-value-labels-1}

for bar, value in zip(bars2, category\_service{[}`avg\_days\_late'{]}):
ax2.text(bar.get\_x() + bar.get\_width()/2, bar.get\_height() + 0.05,
f'\{value:.2f\}`, ha='center', va=`bottom', fontweight=`bold',
fontsize=10)

\section{Plot 3: Late shipment distribution (Top
Right)}\label{plot-3-late-shipment-distribution-top-right}

ax3 = fig.add\_subplot(gs{[}0, 2{]}) late\_by\_category =
shipments\_with\_category{[}shipments\_with\_category{[}`is\_late'{]} ==
True{]} category\_late\_counts =
late\_by\_category{[}`prodCategory'{]}.value\_counts() wedges, texts,
autotexts = ax3.pie(category\_late\_counts.values,
labels=category\_late\_counts.index, colors=colors,
autopct=`\%1.1f\%\%', startangle=90) ax3.set\_title(`Distribution of
Late Shipments\nAcross Categories', fontsize=12, fontweight=`bold',
pad=15)

\section{Plot 4: Multi-category vs Single-category performance (Bottom
Left)}\label{plot-4-multi-category-vs-single-category-performance-bottom-left}

ax4 = fig.add\_subplot(gs{[}1, 0{]}) multi\_cat\_performance =
comprehensive\_analysis.groupby(`has\_multiple\_categories'){[}`is\_late'{]}.mean()
labels = {[}`Single Category', `Multi Category'{]} values =
{[}multi\_cat\_performance{[}False{]} * 100,
multi\_cat\_performance{[}True{]} * 100{]} bars4 = ax4.bar(labels,
values, color={[}`\#4682B4', `\#CD5C5C'{]}, alpha=0.8,
edgecolor=`black') ax4.set\_title(`Late Rate: Single vs Multi-Category
Orders', fontsize=12, fontweight=`bold', pad=15) ax4.set\_ylabel(`Late
Rate (\%)', fontweight=`bold') ax4.set\_ylim(0, max(values) * 1.1)

\section{Add value labels}\label{add-value-labels-2}

for bar, value in zip(bars4, values): ax4.text(bar.get\_x() +
bar.get\_width()/2, bar.get\_height() + 0.5, f'\{value:.1f\}\%`,
ha='center', va=`bottom', fontweight=`bold', fontsize=10)

\section{Plot 5: Service level timeline (Bottom
Center)}\label{plot-5-service-level-timeline-bottom-center}

ax5 = fig.add\_subplot(gs{[}1, 1{]}) \# Create monthly service level
data shipments\_with\_category{[}`month'{]} =
shipments\_with\_category{[}`actualShipDate'{]}.dt.to\_period(`M')
monthly\_service =
shipments\_with\_category.groupby(`month'){[}`is\_late'{]}.mean()
monthly\_on\_time = (1 - monthly\_service) * 100

ax5.plot(range(len(monthly\_on\_time)), monthly\_on\_time.values,
marker=`o', linewidth=2, markersize=6, color=`\#2E8B57')
ax5.set\_title(`Monthly On-Time Delivery Trend', fontsize=12,
fontweight=`bold', pad=15) ax5.set\_ylabel(`On-Time Rate (\%)',
fontweight=`bold') ax5.set\_xlabel(`Month', fontweight=`bold')
ax5.grid(True, alpha=0.3) ax5.set\_ylim(80, 100)

\section{Plot 6: Key metrics summary (Bottom
Right)}\label{plot-6-key-metrics-summary-bottom-right}

ax6 = fig.add\_subplot(gs{[}1, 2{]}) ax6.axis(`off')

\section{Create summary text}\label{create-summary-text}

summary\_text = f''\,``\,'' KEY BUSINESS METRICS

Overall Performance: • On-Time Rate: \{on\_time\_rate:.1f\}\% • Total
Shipments: \{len(clean\_shipments):,\}

Category Performance: • Best: Liquids
(\{category\_service.loc{[}`Liquids', `on\_time\_rate'{]}:.1f\}\%) •
Worst: Machines (\{category\_service.loc{[}`Machines',
`on\_time\_rate'{]}:.1f\}\%)

Multi-Category Orders: • Percentage:
\{multi\_category\_shipments{[}`shipID'{]}.nunique() /
comprehensive\_analysis{[}`shipID'{]}.nunique() * 100:.1f\}\% • Count:
\{multi\_category\_shipments{[}`shipID'{]}.nunique():,\} shipments

Operational Insights: • \{category\_service.loc{[}`Liquids',
`total\_shipments'{]}:,\} Liquids shipments •
\{category\_service.loc{[}`Machines', `total\_shipments'{]}:,\} Machines
shipments • Average delay:
\{category\_service{[}`avg\_days\_late'{]}.mean():.1f\} days ``\,``\,''

ax6.text(0.05, 0.95, summary\_text, transform=ax6.transAxes,
fontsize=10, verticalalignment=`top', fontfamily=`monospace',
bbox=dict(boxstyle=``round,pad=0.5'', facecolor=`lightgray', alpha=0.8))

plt.tight\_layout() plt.show()

\section{Print key insights}\label{print-key-insights}

print(``\n'' + ``=''\emph{60) print(``KEY BUSINESS INSIGHTS'')
print(``=''}60) print(f''• Best performing category: Liquids
(\{category\_service.loc{[}`Liquids', `on\_time\_rate'{]}:.1f\}\%
on-time)``) print(f''• Worst performing category: Machines \& SpareParts
(\{category\_service.loc{[}`Machines', `on\_time\_rate'{]}:.1f\}\%
on-time)``) print(f''• Multi-category orders:
\{multi\_category\_shipments{[}`shipID'{]}.nunique() /
comprehensive\_analysis{[}`shipID'{]}.nunique() * 100:.1f\}\% of
shipments'') print(f''• Overall service level: \{on\_time\_rate:.1f\}\%
on-time delivery'') print(``Visualization code temporarily commented out
for testing'') ``` \#\# 📊 Professional Visualization Analysis: Service
Level Performance by Category

\subsubsection{\texorpdfstring{\textbf{Step 1: Data Foundation and
Business
Context}}{Step 1: Data Foundation and Business Context}}\label{step-1-data-foundation-and-business-context}

\textbf{Dataset Purpose:} This analysis uses ZappTech's shipment data to
understand service level performance across product categories and
identify cross-category ordering patterns.

\textbf{Core Business Questions:} - How does on-time delivery
performance vary across different product categories? - What percentage
of orders include products from multiple categories? - Are there
operational patterns that could inform business strategy?

\textbf{Data Structure:} The dataset combines shipment tracking data
(planned vs actual ship dates) with product category information,
enabling category-level performance analysis.

\begin{center}\rule{0.5\linewidth}{0.5pt}\end{center}

\subsubsection{\texorpdfstring{\textbf{Step 2: Data Quality and
Preparation}}{Step 2: Data Quality and Preparation}}\label{step-2-data-quality-and-preparation}

\textbf{Data Cleaning Process:} - ✅ Removed missing date values to
ensure accurate lateness calculations - ✅ Maintained data integrity by
using left joins to preserve all shipment records - ✅ Calculated
precise lateness metrics using datetime operations

\textbf{Key Metrics Created:} - \texttt{is\_late}: Boolean flag
indicating whether shipment missed planned date - \texttt{days\_late}:
Quantitative measure of delivery delay -
\texttt{has\_multiple\_categories}: Identifies complex multi-category
orders

\begin{center}\rule{0.5\linewidth}{0.5pt}\end{center}

\subsubsection{\texorpdfstring{\textbf{Step 3: Service Level Performance
Analysis}}{Step 3: Service Level Performance Analysis}}\label{step-3-service-level-performance-analysis}

\textbf{Overall Performance Baseline:} The analysis establishes that
\textbf{88.6\%} of shipments deliver on time, providing a benchmark for
category-level comparisons.

\textbf{Category Performance Spectrum:}

\begin{longtable}[]{@{}lll@{}}
\toprule\noalign{}
Category & On-Time Rate & Performance Level \\
\midrule\noalign{}
\endhead
\bottomrule\noalign{}
\endlastfoot
\textbf{Liquids} & \textbf{95.9\%} & 🟢 Excellent \\
\textbf{Marketables} & \textbf{92.2\%} & 🟡 Good \\
\textbf{Machines} & \textbf{82.0\%} & 🔴 Needs Improvement \\
\textbf{Spare Parts} & \textbf{82.1\%} & 🔴 Needs Improvement \\
\end{longtable}

\textbf{Performance Gap:} \textbf{13.9 percentage points} separate best
and worst performing categories, indicating significant operational
variation.

\begin{center}\rule{0.5\linewidth}{0.5pt}\end{center}

\subsubsection{\texorpdfstring{\textbf{Step 4: Multi-Category Order
Analysis}}{Step 4: Multi-Category Order Analysis}}\label{step-4-multi-category-order-analysis}

\textbf{Cross-Customer Behavior:} \textbf{23.3\%} of shipments contain
products from multiple categories, revealing important customer
purchasing patterns.

\textbf{Operational Complexity:} Multi-category orders represent
logistical challenges that may require specialized handling processes.

\textbf{Business Opportunity:} This segment represents potential for
bundled product offerings and targeted marketing strategies.

\begin{center}\rule{0.5\linewidth}{0.5pt}\end{center}

\subsubsection{\texorpdfstring{\textbf{Step 5: Temporal Performance
Trends}}{Step 5: Temporal Performance Trends}}\label{step-5-temporal-performance-trends}

\textbf{Consistency Monitoring:} Monthly trend analysis shows service
level stability over time, helping identify seasonal patterns or
operational changes.

\textbf{Performance Tracking:} The timeline visualization provides
ongoing monitoring capability for service level improvements or
deteriorations.

\begin{center}\rule{0.5\linewidth}{0.5pt}\end{center}

\subsubsection{\texorpdfstring{\textbf{Step 6: Comparative Analysis
Framework}}{Step 6: Comparative Analysis Framework}}\label{step-6-comparative-analysis-framework}

\textbf{Single vs Multi-Category Performance:} Direct comparison reveals
whether complex orders experience different service levels, informing
resource allocation decisions.

\textbf{Category Contribution to Late Shipments:} The distribution
analysis shows which categories contribute most to overall late
shipments, prioritizing improvement efforts.

\begin{center}\rule{0.5\linewidth}{0.5pt}\end{center}

\subsubsection{\texorpdfstring{\textbf{Step 7: Business Intelligence
Integration}}{Step 7: Business Intelligence Integration}}\label{step-7-business-intelligence-integration}

\paragraph{\texorpdfstring{\textbf{Strategic
Insights:}}{Strategic Insights:}}\label{strategic-insights}

\begin{itemize}
\tightlist
\item
  🎯 \textbf{Liquids category} demonstrates operational excellence that
  could be replicated
\item
  ⚠️ \textbf{Machines category} requires focused improvement initiatives
\item
  🔄 \textbf{Multi-category orders} represent both challenge and
  opportunity
\end{itemize}

\paragraph{\texorpdfstring{\textbf{Operational
Recommendations:}}{Operational Recommendations:}}\label{operational-recommendations}

\begin{itemize}
\tightlist
\item
  🔍 Investigate root causes for Machines category performance issues
\item
  📦 Develop specialized processes for multi-category order fulfillment
\item
  📈 Use Liquids category as benchmark for best practices
\end{itemize}

\paragraph{\texorpdfstring{\textbf{Performance
Management:}}{Performance Management:}}\label{performance-management}

\begin{itemize}
\tightlist
\item
  🎯 Establish category-specific service level targets
\item
  📊 Implement regular monitoring of multi-category order trends
\item
  📅 Use monthly trends for proactive operational adjustments
\end{itemize}

\begin{center}\rule{0.5\linewidth}{0.5pt}\end{center}

\subsection{\texorpdfstring{📈 \textbf{Visualization
Components}}{📈 Visualization Components}}\label{visualization-components}

\textbf{Professional Charts Created:} - \textbf{Bar Chart:} Clear
comparison of on-time delivery rates across categories - \textbf{Pie
Chart:} Distribution of late shipments by category\\
- \textbf{Timeline:} Monthly service level trends - \textbf{Summary
Metrics:} Key performance indicators and business insights

\textbf{Key Finding:} The visualization highlights that \textbf{Machines
and SpareParts need operational improvement} while \textbf{Liquids
perform exceptionally well}, providing clear direction for operational
focus areas.

\textbf{Your visualizations should:} - Use clear labels and professional
formatting - Support the insights from your discussion questions - Be
appropriate for a business audience - Do not \texttt{echo} the code that
creates the visualizations

\subsection{Challenge Requirements 📋}\label{challenge-requirements}

\textbf{Your Primary Task:} Answer all discussion questions for the
seven mental models with thoughtful, well-reasoned responses that
demonstrate your understanding of data manipulation concepts.

\textbf{Key Requirements:} - Complete discussion questions for each
mental model - Demonstrate clear understanding of pandas concepts and
data manipulation techniques - Write clear, business-focused analysis
that explains your findings

\subsection{Getting Started: Repository Setup
🚀}\label{getting-started-repository-setup}

\begin{tcolorbox}[enhanced jigsaw, leftrule=.75mm, bottomtitle=1mm, toprule=.15mm, toptitle=1mm, left=2mm, colbacktitle=quarto-callout-important-color!10!white, rightrule=.15mm, titlerule=0mm, colback=white, colframe=quarto-callout-important-color-frame, title=\textcolor{quarto-callout-important-color}{\faExclamation}\hspace{0.5em}{📁 Getting Started}, arc=.35mm, coltitle=black, opacitybacktitle=0.6, breakable, opacityback=0, bottomrule=.15mm]

\textbf{Step 1:} Fork and clone this challenge repository - Go to the
course repository and find the ``dataManipulationChallenge'' folder -
Fork it to your GitHub account, or clone it directly - Open the cloned
repository in Cursor

\textbf{Step 2:} Set up your Python environment - Follow the Python
setup instructions above (use your existing venv from Tech Setup
Challenge Part 2) - Make sure your virtual environment is activated and
the Python interpreter is set

\textbf{Step 3:} You're ready to start! The data loading code is already
provided in this file.

\textbf{Note:} This challenge uses the same \texttt{index.qmd} file
you're reading right now - you'll edit it to complete your analysis.

\end{tcolorbox}

\subsubsection{Getting Started Tips}\label{getting-started-tips}

\begin{tcolorbox}[enhanced jigsaw, leftrule=.75mm, bottomtitle=1mm, toprule=.15mm, toptitle=1mm, left=2mm, colbacktitle=quarto-callout-note-color!10!white, rightrule=.15mm, titlerule=0mm, colback=white, colframe=quarto-callout-note-color-frame, title=\textcolor{quarto-callout-note-color}{\faInfo}\hspace{0.5em}{🎯 Method Chaining Philosophy}, arc=.35mm, coltitle=black, opacitybacktitle=0.6, breakable, opacityback=0, bottomrule=.15mm]

\begin{quote}
``Each operation should build naturally on the previous one''
\end{quote}

\emph{Think of method chaining like building with LEGO blocks - each
piece connects to the next, creating something more complex and useful
than the individual pieces.}

\end{tcolorbox}

\begin{tcolorbox}[enhanced jigsaw, leftrule=.75mm, bottomtitle=1mm, toprule=.15mm, toptitle=1mm, left=2mm, colbacktitle=quarto-callout-warning-color!10!white, rightrule=.15mm, titlerule=0mm, colback=white, colframe=quarto-callout-warning-color-frame, title=\textcolor{quarto-callout-warning-color}{\faExclamationTriangle}\hspace{0.5em}{💾 Important: Save Your Work Frequently!}, arc=.35mm, coltitle=black, opacitybacktitle=0.6, breakable, opacityback=0, bottomrule=.15mm]

\textbf{Before you start:} Make sure to commit your work often using the
Source Control panel in Cursor (Ctrl+Shift+G or Cmd+Shift+G). This
prevents the AI from overwriting your progress and ensures you don't
lose your work.

\textbf{Commit after each major step:}

\begin{itemize}
\tightlist
\item
  After completing each mental model section
\item
  After adding your visualizations
\item
  After completing your advanced method chain
\item
  Before asking the AI for help with new code
\end{itemize}

\textbf{How to commit:}

\begin{enumerate}
\def\labelenumi{\arabic{enumi}.}
\tightlist
\item
  Open Source Control panel (Ctrl+Shift+G)
\item
  Stage your changes (+ button)
\item
  Write a descriptive commit message
\item
  Click the checkmark to commit
\end{enumerate}

\emph{Remember: Frequent commits are your safety net!}

\end{tcolorbox}

\subsection{Grading Rubric 🎓}\label{grading-rubric}

\textbf{75\% Grade:} Complete discussion questions for at least 5 of the
7 mental models with clear, thoughtful responses.

\textbf{85\% Grade:} Complete discussion questions for all 7 mental
models with comprehensive, well-reasoned responses.

\textbf{95\% Grade:} Complete all discussion questions plus the
``Answering A Business Question'' section.

\textbf{100\% Grade:} Complete all discussion questions plus create a
professional visualization showing service level by product category.

\subsection{Submission Checklist ✅}\label{submission-checklist}

\textbf{Minimum Requirements (Required for Any Points):}

\begin{itemize}
\tightlist
\item[$\square$]
  Created repository named ``dataManipulationChallenge'' in your GitHub
  account
\item[$\square$]
  Cloned repository locally using Cursor (or VS Code)
\item[$\square$]
  Completed discussion questions for at least 5 of the 7 mental models
\item[$\square$]
  Document rendered to HTML successfully
\item[$\square$]
  HTML files uploaded to your repository
\item[$\square$]
  GitHub Pages enabled and working
\item[$\square$]
  Site accessible at
  \texttt{https://{[}your-username{]}.github.io/dataManipulationChallenge/}
\end{itemize}

\textbf{75\% Grade Requirements:}

\begin{itemize}
\tightlist
\item[$\square$]
  Complete discussion questions for at least 5 of the 7 mental models
\item[$\square$]
  Clear, thoughtful responses that demonstrate understanding
\item[$\square$]
  Code examples where appropriate to support explanations
\end{itemize}

\textbf{85\% Grade Requirements:}

\begin{itemize}
\tightlist
\item[$\square$]
  Complete discussion questions for all 7 mental models
\item[$\square$]
  Comprehensive, well-reasoned responses showing deep understanding
\item[$\square$]
  Business context for why concepts matter
\item[$\square$]
  Examples of real-world applications
\end{itemize}

\textbf{95\% Grade Requirements:}

\begin{itemize}
\tightlist
\item[$\square$]
  Complete discussion questions for all 7 mental models
\item[$\square$]
  Complete the ``Answering A Business Question'' discussion questions
\item[$\square$]
  Comprehensive, well-reasoned responses showing deep understanding
\item[$\square$]
  Business context for why concepts matter
\end{itemize}

\textbf{100\% Grade Requirements:}

\begin{itemize}
\tightlist
\item[$\square$]
  All discussion questions completed with professional quality
\item[$\square$]
  Professional visualization showing service level by product category
\item[$\square$]
  Professional presentation style appropriate for business audience
\item[$\square$]
  Clear, engaging narrative that tells a compelling story
\item[$\square$]
  Practical insights that would help ZappTech's management
\end{itemize}

\textbf{Report Quality (Critical for Higher Grades):}

\begin{itemize}
\tightlist
\item[$\square$]
  Professional writing style (no AI-generated fluff)
\item[$\square$]
  Concise analysis that gets to the point
\end{itemize}




\end{document}
